\documentclass[12pt,oneside]{book}

% ==================== METADATOS ====================
\newcommand{\universidad}{Universidad de Buenos Aires (UBA)}
\newcommand{\facultad}{Facultad de Derecho}
\newcommand{\curso}{Curso de Perfeccionamiento --- Código Civil y Comercial 2015 (Parte 2)}
\newcommand{\materia}{Derecho de las Sucesiones}
\newcommand{\titulo}{Sucesión intestada, sucesión del cónyuge y legítima}
\newcommand{\autor}{Isaac Edgar Camacho Ocampo}
\newcommand{\anio}{2015}

% ==================== PAQUETES ====================
\usepackage[utf8]{inputenc}
\usepackage[T1]{fontenc}
\usepackage[spanish,provide=*]{babel}
\usepackage[
    top=2.5cm,
    bottom=2.5cm,
    left=3cm,
    right=2.5cm,
    headheight=15pt
]{geometry}
\usepackage{amsmath}
\usepackage{amssymb}
\usepackage{mathtools}
\usepackage{graphicx}
\usepackage{float}
\usepackage{caption}
\usepackage{subcaption}
\usepackage{booktabs}
\usepackage{array}
\usepackage{longtable}
\usepackage{tabularx}
\usepackage{enumitem}
\usepackage{xcolor}
\usepackage[most]{tcolorbox}
\usepackage{fancyhdr}
\usepackage{ragged2e}
\usepackage[
    colorlinks=true,
    linkcolor=blue!50!black,
    citecolor=green!40!black,
    urlcolor=blue!60!black,
    pdftitle={\titulo},
    pdfauthor={\autor},
    pdfsubject={Apuntes universitarios},
    bookmarks=true
]{hyperref}

% ==================== CONFIGURACIÓN ====================
\graphicspath{
    {./img/}
    {./graficos/}
    {./figuras/}
}
\setcounter{tocdepth}{2}
\setlength{\parindent}{0pt}
\setlength{\parskip}{0.5em}
\setlist[itemize]{
    topsep=0.3em,
    itemsep=0.2em
}
\setlist[enumerate]{
    topsep=0.3em,
    itemsep=0.2em
}
\clubpenalty=10000
\widowpenalty=10000

% ==================== COMANDOS ====================
\newcommand{\abs}[1]{\left|#1\right|}
\newcommand{\norm}[1]{\left\lVert#1\right\rVert}
\newcommand{\vect}[1]{\mathbf{#1}}
\newcommand{\deriv}[2]{\frac{d#1}{d#2}}
\newcommand{\pderiv}[2]{\frac{\partial#1}{\partial#2}}

% ==================== ENTORNOS / CAJAS ====================
\newtcolorbox{definition}[1]{
    enhanced,
    breakable,
    title={Definición: #1},
    fonttitle=\bfseries,
    colback=blue!3,
    colframe=blue!50!black,
    boxrule=0.8pt,
    arc=2mm
}

\newtcolorbox{important}{
    enhanced,
    breakable,
    title={Importante},
    fonttitle=\bfseries,
    colback=yellow!8,
    colframe=orange!70!black,
    boxrule=0.8pt,
    arc=2mm
}

\newtcolorbox{example}[1]{
    enhanced,
    breakable,
    title={Ejemplo: #1},
    fonttitle=\bfseries,
    colback=green!4,
    colframe=green!40!black,
    boxrule=0.8pt,
    arc=2mm
}

\newtcolorbox{formula}[1]{
    enhanced,
    breakable,
    title={Fórmula / Regla: #1},
    fonttitle=\bfseries,
    colback=purple!4,
    colframe=purple!50!black,
    boxrule=0.8pt,
    arc=2mm
}

\newtcolorbox{exercise}[1]{
    enhanced,
    breakable,
    title={Ejercicio: #1},
    fonttitle=\bfseries,
    colback=gray!5,
    colframe=gray!60!black,
    boxrule=0.8pt,
    arc=2mm
}

\newtcolorbox{note}{
    enhanced,
    breakable,
    title={Nota},
    fonttitle=\bfseries,
    colback=white,
    colframe=gray!50,
    boxrule=0.6pt,
    arc=2mm
}

% ==================== CONFIGURACIÓN FINAL ====================
\pagestyle{fancy}
\fancyhf{}
\fancyhead[L]{\nouppercase{\leftmark}}
\fancyhead[R]{\materia}
\fancyfoot[C]{\thepage}
\renewcommand{\headrulewidth}{0.4pt}
\renewcommand{\footrulewidth}{0pt}

\fancypagestyle{plain}{
    \fancyhf{}
    \fancyfoot[C]{\thepage}
    \renewcommand{\headrulewidth}{0pt}
}

\renewcommand{\arraystretch}{1.25}

\begin{document}

% ==================== PORTADA ====================
\begin{titlepage}

\thispagestyle{empty}

\begin{center}

\vspace*{1cm}

{\large \textsc{\universidad}}

\vspace{0.2cm}

{\large \textsc{\facultad}}

\vspace{0.2cm}

{\large \curso}

\vspace{3cm}

{\Huge \textbf{\materia}}

\vspace{1cm}

{\LARGE \titulo}

\vspace{0.8cm}

\textit{Comprender cada concepto es el primer paso para dominarlo.}

\vspace{1.2cm}

\rule{0.70\textwidth}{0.5pt}

\vspace{0.5cm}

{\large Apuntes de estudio}

\vspace{0.5cm}

\rule{0.70\textwidth}{0.5pt}

\vfill

{\large \autor}

\vspace{0.3cm}

{\large \anio}

\end{center}

\vfill

\begin{flushleft}
\textit{Este es un modesto aporte a los estudiantes de ``Derecho de las Sucesiones'' no pretende sustituir el material oficial de la materia.}
\end{flushleft}

\end{titlepage}

% ==================== ÍNDICE ====================
\tableofcontents
\clearpage

% ==================== CONTENIDO ====================

\chapter*{Presentación de la clase}
\addcontentsline{toc}{chapter}{Presentación de la clase}
\markboth{Presentación de la clase}{}

\section*{Temas tratados}

La clase compara el régimen sucesorio del Código de Vélez con el del Código Civil y Comercial (CCyC). En primer lugar recorre la \textbf{sucesión intestada} ---fuentes del llamamiento, órdenes, grados, derecho de representación y herencia de descendientes, ascendientes, adoptados y colaterales--- y detalla, con mayor profundidad, la \textbf{sucesión del cónyuge}: separación de bienes, sociedad conyugal disuelta, cónyuge putativo, habitación viudal y causales de exclusión (matrimonio \textit{in extremis}, divorcio y separación de hecho).

La segunda mitad estudia la \textbf{legítima}: cuotas, cálculo, protección, acción de reducción, sus efectos frente a donatarios y terceros, y la transmisión de bienes a legitimarios con reserva de usufructo.

\begin{note}
La clase se dictó antes de la entrada en vigencia del Código Civil y Comercial (1.\textordmasculine{} de agosto de 2015); las referencias al «Código actual» corresponden al Código de Vélez. De los artículos citados, solo el art. 2445 se transcribe del texto oficial; los demás se consignan según lo leído o expuesto en clase, por lo que conviene cotejarlos con el texto vigente.
\end{note}
% TODO: verificar en la fuente: el número «art. 714» (caducidad de la acción de nulidad por muerte, equivalente al art. 239 de Vélez); los casos «Escary c/ Pietranera», plenario «Mauri», plenario «Espota», «Bogelius», «Astesiano» y «Morro Berner»; y el nombre de la posición doctrinaria «de Fuentes».

\section*{Relevancia profesional y personal}

Casi todo caso de familia, herencia o planificación patrimonial termina pasando por estas reglas. Quien litiga necesita saber quién hereda, en qué proporción, qué parte puede disponer el causante y cómo se ataca una donación que perjudica la legítima; quien ejerce como notario debe poder asesorar sobre testamentos, donaciones con reserva de usufructo y pactos de herencia futura sin exponer a su cliente a una reducción. Además, la clase enseña a leer una reforma con ojo crítico: detectar contradicciones (arts. 2459 y 2461), silencios legales y puntos que la jurisprudencia deberá resolver.

El contenido es útil más allá del estudio: todas las personas forman una familia, tienen bienes y alguna vez les tocará heredar o planificar una herencia. Entender por qué una persona no puede dejar todo a quien quiere, qué pasa con la pareja que convive pero no se casó, o por qué conviene formalizar una donación, permite tomar decisiones informadas y evitar conflictos familiares, porque el derecho sucesorio decide quién recibe lo que una vida entera construyó.

\section*{Cómo se relacionan los temas: un ejemplo cotidiano}

\begin{example}{Ricardo y su herencia}
Ricardo fallece. Estuvo casado con Marta ---separados de hecho hace años---, tiene dos hijos de ese matrimonio, convive desde hace un año con Laura y donó a su amigo Tito, hace doce años, la casa de la costa.

\begin{itemize}
    \item \textbf{Capítulo 1:} Laura pregunta si su convivencia le da algún derecho.
    \item \textbf{Capítulo 2:} se ordena quién hereda: los hijos por derecho propio y, si alguno hubiera muerto, sus hijos por representación.
    \item \textbf{Capítulo 3:} se analiza si Marta hereda o queda excluida por la separación de hecho.
    \item \textbf{Capítulo 4:} se calcula cuánta parte de lo dejado es legítima de los hijos y cuánta podía disponer Ricardo.
    \item \textbf{Capítulo 5:} se ve qué ocurre con la casa donada a Tito: si perjudicó la legítima, los hijos pueden iniciar la acción de reducción.
\end{itemize}
\end{example}

% ------------------------------------------------------------------
\chapter[Consultas previas de los alumnos]{Consultas previas de los alumnos}

Este capítulo reúne las consultas formuladas por los alumnos sobre las unidades 1 a 3 y las respuestas de la expositora de la primera parte (sucesión intestada).

\section{La prueba del testamento cuando no hay pericia caligráfica}

\textit{Consulta: \textbf{\textquestiondown{}Cómo se prueba un testamento cuando no es posible una pericia caligráfica?}}

En el Código actual se exigen dos testigos. Si no hay testigos, se recurre a la pericia caligráfica, que cuenta con jurisprudencia que la ha aceptado cuando faltan testigos. Surge entonces la pregunta por el testamento respecto del cual no existe posibilidad de pericia caligráfica: debe poder probarse por cualquier otro medio, y la expositora entiende que ese medio serían los \textbf{testigos}.

No puede dejar de protocolizarse un testamento porque la prueba no esté efectivamente establecida; lo que corresponde, si no puede realizarse la pericia caligráfica, es requerir otra prueba. La expositora aclara que se trata de una opinión propia, formulada como juez, y que todo ello es opinable mientras no exista una posición jurisprudencial definida.

\section{Hija única y separación de bienes}

\begin{quote}
\textit{Consulta: «El causante es propietario del 50\,\% del inmueble y tiene solamente una hija. Con el nuevo Código, \textquestiondown{}varía en algo lo que sería una sucesión actual?»}
\end{quote}

\textbf{No varía.} En materia de sucesión intestada hay pocas variaciones en cuanto a la sucesión de los descendientes y del cónyuge.

\begin{quote}
\textit{Consulta: «En el régimen de separación de bienes, \textquestiondown{}el otro cónyuge hereda como un hijo más, o se disuelve la sociedad conyugal y se le asignan igualmente los bienes gananciales?»}
\end{quote}

La pregunta encierra un contrasentido: si hay separación de bienes, \textbf{no hay bienes gananciales}. El CCyC no establece cómo hereda el cónyuge en caso de separación de bienes, que no son ni bienes propios ni gananciales, sino \textbf{bienes personales} de cada cónyuge.

Hasta ahora se aplica \textbf{por analogía} la solución prevista para los bienes propios: en la separación de bienes el cónyuge hereda con los hijos \textbf{por cabeza}, es decir, como un hijo más. Esto debería haberse establecido expresamente, porque el art. 508, al ocuparse de la separación de bienes, se remite a la división por la metodología de las herencias, lo que puede generar discordancias. Pero ya no se trata de liquidar una sociedad conyugal inexistente, sino de determinar la propiedad de los bienes cuando no pueda establecerse de quién es alguno (el Código lo resuelve como común) o cuando exista algún bien en condominio, que se dividirá en el ámbito de la partición de herencias según la remisión del art. 508.

\section{El concubino y la vocación hereditaria}

\begin{quote}
\textit{Consulta: «En el nuevo Código, \textquestiondown{}se le otorga algún tipo de vocación hereditaria al concubino?»}
\end{quote}

\textbf{No.} Existen autores que critican al Código por no haber otorgado derecho sucesorio al concubino: al regularse las uniones convivenciales, podía pensarse que se reconocería algún derecho hereditario tras determinados años de convivencia, pero no se lo hizo. En este aspecto \textbf{todo sigue igual}.

El concubino no tiene ningún derecho hereditario en el CCyC, aunque podría tener más derechos que en el matrimonio.

% ------------------------------------------------------------------
\chapter[Sucesión intestada]{Sucesión intestada: lo que se mantiene}

Este capítulo desarrolla la sucesión intestada, con las modificaciones introducidas por el CCyC, algunas favorables y otras no. La expositora es la docente de la primera parte.

\section{Fuentes del llamamiento, apertura y transmisión de la herencia}

\subsection{Las fuentes del llamamiento y el momento de la transmisión}

En el CCyC no varían las fuentes del llamamiento: la que deriva de la \textbf{ley} y la que deriva del \textbf{testamento}. Tampoco cambió que \textbf{la muerte, la apertura de la sucesión y la transmisión de la herencia se producen en un mismo momento}. Esta concepción, que en el Código de Vélez surgía de una nota, se encuentra ahora en el texto del \textbf{art. 2277}, de modo que se mantiene la misma concepción.

También se conserva la concepción velezana según la cual se puede \textbf{morir en parte testado y en parte intestado}: la sucesión intestada se abre aun cuando haya testamento, si en este no se instituyen herederos o, instituyéndose herederos, existen legitimarios.

\subsection{De «heredero forzoso» a «legitimario»}

El CCyC ya no utiliza la terminología \textbf{heredero forzoso}, que venía de Roma, pasando por el código español, hasta llegar a Vélez. Ahora se habla directamente de \textbf{legitimario}. A criterio de la expositora, esto refuerza la tendencia a la legítima \textit{pars bonorum} en el nuevo código, cuestión que desarrollará el docente a cargo de la legítima.

\subsection{Casos en que se abre la sucesión intestada}

La sucesión intestada se abre en los mismos casos que en el Código de Vélez:

\begin{enumerate}
    \item cuando el testamento no instituye herederos (por ejemplo, instituye solo legados y queda un remanente);
    \item cuando no hay testamento;
    \item cuando el testamento falla, por nulidad o revocación;
    \item cuando hay legitimarios a quienes se otorga menos de la legítima o se los priva de ella.
\end{enumerate}

\section{Órdenes, grados y origen de los bienes}

\subsection{Los órdenes hereditarios}

Se mantiene la transmisión por órdenes y se mantienen los mismos órdenes: los descendientes como primer orden; los ascendientes y el cónyuge, que constituyen un orden «anómalo»; y los colaterales como tercer orden. El cónyuge excluye a los colaterales. En cuanto a la categoría de órdenes, nada ha variado.

\begin{table}[H]
\centering
\small
\begin{tabularx}{\textwidth}{
    >{\raggedright\arraybackslash}p{1.2cm}
    >{\raggedright\arraybackslash}p{3.6cm}
    >{\raggedright\arraybackslash}X
}
\toprule
\textbf{Orden} & \textbf{Quiénes} & \textbf{Características que se mantienen} \\
\midrule
1.\textordmasculine{} & Descendientes & Heredan por derecho propio los de primer grado; los restantes, por representación. Legítima de dos tercios (art. 2445) e investidura de pleno derecho (art. 2337). \\
\addlinespace
2.\textordmasculine{} & Ascendientes y cónyuge (orden «anómalo») & Los ascendientes tienen legítima, investidura de pleno derecho y no son sujetos pasivos de la colación. El cónyuge excluye a los colaterales. \\
\addlinespace
3.\textordmasculine{} & Colaterales & Hasta el cuarto grado. No tienen investidura de pleno derecho (deben pedirla a los jueces) ni legítima. Solo hay representación en los descendientes de hermanos. \\
\bottomrule
\end{tabularx}
\caption*{Los órdenes hereditarios en la sucesión intestada.}
\end{table}

\subsection{Los grados}

Dentro de cada orden se hereda por grados, salvo el derecho de representación. Tampoco han variado los grados: en general los grados más cercanos excluyen a los más remotos, salvo en el derecho de representación, que opera en dos líneas: la línea descendente y la primera línea colateral, es decir, los descendientes de los hermanos. Se explicará más adelante en qué aspecto sí ha variado el derecho de representación.

\subsection{El origen de los bienes y sus excepciones}

El CCyC tampoco modificó el principio de que \textbf{no se tiene en cuenta el origen de los bienes}, aunque la expositora considera que en el nuevo código pueden encontrarse más excepciones. En el Código actual hay dos excepciones, que subsisten:

\begin{enumerate}
    \item \textbf{Bienes propios y gananciales cuando hereda el cónyuge}: cuando el fallecimiento se produce durante la vigencia de la sociedad conyugal, o con una sociedad conyugal disuelta pero no liquidada, se diferencian los bienes propios y gananciales y se determina la calificación de los bienes.
    \item \textbf{Sucesión del adoptado}: el adoptante no hereda los bienes que a título gratuito hubiera recibido el adoptado de la familia de origen, y la familia de origen no hereda los bienes que a título gratuito hubiera recibido el adoptado de los adoptantes. El Código de Vélez tenía una falencia ---la \textbf{vacancia}, cuando había bienes a título gratuito y no existían familia de origen o familia adoptante a la muerte del adoptado---, que ahora está subsanada: los bienes corresponden a uno u otro según el caso.
\end{enumerate}

A ellas se agregan otras dos diferenciaciones de bienes dentro de la sucesión:

\paragraph{Derecho de opción (art. 2290).} El causante pudo morir sin haber aceptado ni repudiado una herencia: el derecho de opción entre aceptarla o repudiarla integra su patrimonio hereditario. Si algunos de sus herederos aceptan y otros renuncian, en la sucesión habrá que distinguir los bienes de quienes aceptaron de los de quienes no.

\paragraph{Aceptación forzada (art. 2295).} Quien acepta forzadamente (por haber ocultado bienes) no participa de esos bienes: serán heredados por los demás herederos, mientras que el aceptante forzoso no hereda los bienes que ocultó. También aquí habrá que diferenciar bienes dentro de la sucesión.

\section{El derecho de representación}

No hay diferencias respecto del Código de Vélez en cuanto a quién se representa: se representa \textbf{al muerto, al indigno y al renunciante}. Ya no se representa al \textbf{desheredado}, porque \textbf{no existe más la desheredación}. Aparecen dos temas en materia de representación.

\subsection{Sucesión por estirpe}

En el derecho de representación se sucede por \textbf{estirpe}: se diferencian las distintas ramas y, dentro de cada rama, se hereda por cabeza o en partes iguales. Cuando el nieto concurre a la sucesión del abuelo con algún tío, la sucesión por estirpe tiene sentido.

\begin{example}{la injusticia de la estirpe entre nietos}
Si todas las líneas intermedias han fallecido, renunciado o sido declaradas indignas y vienen solamente nietos de distintas ramas, todos son nietos del causante. Sin embargo, por la sucesión por estirpe, si los nietos de una rama son cinco y los de otra rama es uno, el reparto resulta muy injusto: lo lógico sería que heredaran en partes iguales.
\end{example}

Esta solución fue criticada en el Código Civil y \textbf{no fue modificada}: continúa vigente en el \textbf{art. 2428}. Los nietos vienen siempre por derecho de representación.

\subsection{Representación en la sucesión testamentaria (art. 2429)}

Se ha resuelto la discusión doctrinaria ---también práctica--- sobre si había derecho de representación en la sucesión testamentaria. Vélez consideraba que, si se hacía una institución a los parientes, se ejercía el derecho de representación. Basados en ese artículo, muchos consideraron que en la sucesión testamentaria podía reconocerse el derecho de representación siempre que el testador reprodujera el llamamiento legal; otra parte de la doctrina entendía que, de haberlo querido, el testador tendría que haber hecho una sustitución permitida, por lo que no habría representación. Se discutía incluso si procedía solo cuando se reproducía la ley o solo para los legitimarios: había tres o cuatro posiciones doctrinarias.

\begin{formula}{Art. 2429 CCyC --- Representación en la sucesión testamentaria}
«También se aplica el derecho de representación en la sucesión testamentaria si el testador se limita a confirmar la distribución de la herencia que resulta de la ley.»
\end{formula}

Es una \textbf{presunción de la ley}: siempre hay derecho de representación si el testador reprodujo el llamamiento legal. Con ello desaparece la discusión. En el derecho de representación no hay ninguna otra modificación.

\section[Descendientes, ascendientes y colaterales]{Descendientes, ascendientes, adoptantes y colaterales}

\subsection{Descendientes y adoptados}

Al igual que en el código actual, los únicos que heredan por derecho propio son los \textbf{descendientes de primer grado}. Los restantes (nietos, bisnietos, etc.) heredan siempre por derecho de representación.

Hay una modificación respecto del adoptado: el CCyC \textbf{no diferencia entre adopción simple y adopción plena}, de modo que todo adoptado hereda como si fuera un hijo (natural o por técnicas de reproducción asistida). Esa diferenciación subsiste en el caso del adoptante y respecto de los ascendientes, pero no en el de los descendientes: \textbf{el adoptado hereda como un hijo más}. Antes había diferenciación en el caso de la adopción simple.

El descendiente es legitimario, con una legítima de \textbf{dos tercios} (art. 2445) y tiene la \textbf{investidura de pleno derecho} (art. 2337). En algunos casos excluye al cónyuge en los mismos casos que actualmente, es decir, cuando en la herencia solo hay bienes gananciales. En lo que se aplica todos los días se mantiene exactamente lo que dice el Código Civil.

\subsection{Ascendientes}

Los ascendientes tienen legítima, que se ha \textbf{disminuido a la mitad}. También tienen la investidura (posesión de pleno derecho) por el art. 2337. No son sujetos pasivos de la colación: el nuevo código considera que una donación a un ascendiente no es un anticipo de herencia, como sí cabe presumirlo de una donación a un descendiente.

El problema es si el ascendiente es \textbf{sujeto activo} de la colación, por ser legitimario. Hay quien sostiene que tendría la acción de colación; otros lo niegan por reciprocidad: si no se lo considera sujeto pasivo, tampoco podría ser sujeto activo. La expositora considera más justa esta segunda posición, aunque todo son conjeturas hasta que haya doctrina y jurisprudencia.

\subsection{Adoptantes y familia de origen (art. 2432)}

El art. 2432 constituye una \textbf{innovación}: ya no quedará vacancia, sino que heredarán los padres de sangre o los padres de adopción, según quien exista, respecto de los bienes adquiridos a título gratuito. El resto de los bienes ---esto sigue igual--- lo hereda siempre el adoptante.

\subsection{Colaterales}

Los colaterales heredan hasta el \textbf{cuarto grado}, igual que ahora. No tienen investidura de pleno derecho (deben pedirla a los jueces) ni legítima. La única representación es la de los descendientes de los hermanos; los demás colaterales, por grado, se excluyen unos a otros y no tienen derecho de representación. También se mantiene la diferenciación entre hermanos y medio hermanos: el \textbf{medio hermano hereda la mitad} que el hermano bilateral. Queda excluida la nuera viuda sin hijos.

\begin{important}
\textbf{Nada cambia para los parientes en la sucesión intestada.}
\begin{itemize}
    \item \textbf{Descendientes:} heredan por derecho propio los de primer grado; los demás, por representación. El adoptado hereda como un hijo más (ya no se distingue adopción simple y plena).
    \item \textbf{Ascendientes:} legítima reducida a la mitad; no son sujetos pasivos de la colación.
    \item \textbf{Adoptantes y familia de origen:} se subsana la vacancia (art. 2432).
    \item \textbf{Colaterales:} hasta el cuarto grado, sin investidura de pleno derecho y sin legítima; el medio hermano hereda la mitad que el bilateral; solo hay representación en los descendientes de hermanos.
\end{itemize}
\end{important}

\section{Cuadro Cornell: sucesión intestada}

\begin{longtable}{
    >{\raggedright\arraybackslash}p{0.27\textwidth}
    >{\raggedright\arraybackslash}p{0.67\textwidth}
}
\toprule
\textbf{Preguntas / palabras clave} &
\textbf{Notas / respuestas} \\
\midrule
\endfirsthead
\toprule
\textbf{Preguntas / palabras clave} &
\textbf{Notas / respuestas} \\
\midrule
\endhead
\bottomrule
\endfoot

\textbf{\textquestiondown{}Cuándo se abre la sucesión intestada en el nuevo Código y cuándo se transmite la herencia?} &
Cuando el testamento no instituye herederos (solo legados con remanente), no hay testamento, falla por nulidad o revocación, o hay legitimarios a quienes se otorga menos de la legítima o se los priva de ella. La muerte, la apertura y la transmisión de la herencia se producen en un mismo momento (art. 2277). \\[0.6em]

\textbf{\textquestiondown{}Qué órdenes y grados se mantienen?} &
Descendientes; ascendientes y cónyuge; colaterales. Dentro de cada orden se hereda por grados (el más cercano excluye al más remoto), salvo el derecho de representación. Se mantiene el desinterés por el origen de los bienes, con excepciones. \\[0.6em]

\textbf{\textquestiondown{}Qué excepciones hay al principio de que no importa el origen de los bienes?} &
Bienes propios y gananciales cuando hereda el cónyuge; sucesión del adoptado (se subsana la vacancia, art. 2432); derecho de opción transmitido (art. 2290); y aceptación forzada (art. 2295). \\[0.6em]

\textbf{\textquestiondown{}Qué novedad trae el derecho de representación?} &
Se representa al muerto, al indigno y al renunciante, no al desheredado (ya no existe la desheredación). El art. 2429 resuelve la discusión: hay representación también en la sucesión testamentaria si el testador solo confirma el llamamiento legal. La sucesión por estirpe (art. 2428) no fue modificada. \\[0.6em]

\textbf{\textquestiondown{}Cómo hereda el adoptado?} &
Como un hijo más: ya no se distingue entre adopción simple y plena en la sucesión de descendientes. Sí se mantienen diferencias respecto del adoptante y de la familia de origen. \\[0.6em]

\midrule
\multicolumn{2}{p{\dimexpr0.94\textwidth+2\tabcolsep\relax}}{
\textbf{Resumen:} El CCyC conserva casi intacta la sucesión intestada del Código de Vélez: fuentes del llamamiento, transmisión por órdenes y grados, derecho de representación y sucesión de descendientes, ascendientes y colaterales. Las mejoras puntuales son la subsanación de la vacancia en la herencia del adoptado y la solución del art. 2429 sobre la representación en la sucesión testamentaria.
} \\
\end{longtable}

% ------------------------------------------------------------------
\chapter[Sucesión del cónyuge]{La sucesión del cónyuge}

\section[Panorama y derecho comparado]{Panorama de la sucesión del cónyuge y derecho comparado}

En la sucesión del cónyuge sí es necesario hacer diferenciaciones de acuerdo con las reformas. Incide en ella el \textbf{divorcio incausado}, que repercute en la separación de hecho sin voluntad de unirse y en el divorcio como causal de exclusión.

La sucesión del cónyuge siempre tuvo problemas: por un lado, con el cónyuge pasaban los bienes de una familia a la otra; por otro, existía la incapacidad de la mujer. La evolución de la mujer de la incapacidad a la capacidad y, sobre todo, la introducción del divorcio en las distintas legislaciones modificó la problemática: ahora puede haber sucesivos hijos de sucesivos matrimonios y un cónyuge. Así, el cónyuge puede ser excluido por descendientes que no son sus propios hijos, por ejemplo cuando en la herencia solo hay bienes gananciales. Por eso hay legislaciones que tratan la sucesión del cónyuge teniendo en cuenta principalmente los \textbf{sucesivos matrimonios}: el divorcio y la posibilidad de sucesivos matrimonios inciden mucho más en la sucesión que en el régimen patrimonial del matrimonio.

El CCyC \textbf{no estuvo a la altura de los tiempos}, porque mantuvo prácticamente igual la sucesión del cónyuge. Otros derechos prevén soluciones distintas, mencionadas a título informativo:

\begin{itemize}
    \item el derecho español hace recaer un usufructo sobre la parte disponible (o una parte del tercio de mejora) cuando el cónyuge concurre con muchos hijos de distintos matrimonios, para no perjudicar la legítima de los hijos;
    \item el código suizo hace optar entre la propiedad y el usufructo: menor porción si se opta por la propiedad y mayor porción si se opta por el usufructo.
\end{itemize}

\section{Separación de bienes y artículo 2433}

\subsection{El artículo 2433 y los arts. 3570 y 3576 de Vélez}

El art. 2433 reemplaza a los arts. 3570 y 3576 del Código de Vélez. El art. 3570 trataba los bienes propios, que el cónyuge heredaba concurriendo por cabeza con los hijos. El art. 3576 trataba la exclusión del cónyuge cuando había descendientes sobre los bienes gananciales, y sufrió varias modificaciones; actualmente todo tipo de descendiente, matrimonial o extramatrimonial, excluye al cónyuge sobre los bienes gananciales. El art. 2433 \textbf{reune ambas reglas}: su primera parte trata los bienes propios (el cónyuge hereda por cabeza, como un hijo más) y la segunda los bienes gananciales (los descendientes excluyen al cónyuge).

El art. 2433 solo prevé la hipótesis de un matrimonio con disolución de la sociedad conyugal por muerte; \textbf{no previó la muerte de un cónyuge durante la vigencia de la separación de bienes}. La expositora considera que esto está mal y que debería haberse diferenciado, al menos en la primera parte, cuando habla de bienes propios: en un régimen de separación de bienes no hay bienes propios, sino bienes personales de cada cónyuge.

\subsection{Crítica al régimen de separación de bienes}

Para la expositora, el régimen de separación de bienes es \textbf{muy injusto para el débil de la relación}, que generalmente es la mujer, porque ni siquiera se ha establecido una compensación. En España, por ejemplo, se establece una compensación en el régimen de separación para quien realiza trabajos en el hogar. Mientras uno cría a los hijos, el otro resulta titular de todos los bienes sin compensación alguna, a diferencia de otras legislaciones europeas: es un régimen muy ``egoísta''.

Para la expositora, el mejor régimen cuando hay convivencia es la \textbf{sociedad conyugal} (comunidad), difundido en todas las legislaciones del mundo; incluso en legislaciones ajenas se permiten y se hacen reiteradamente pactos acogiéndose al régimen de comunidad. Si se quería otra opción, habría correspondido el régimen de participación y no el de separación.

\subsection{El condominio y la partición en la separación de bienes}

En la separación de bienes cada uno es propietario de lo que adquiere; no hay ganancialidad. Lo único que prevé el Código es que no pueda determinarse de quién es un bien adquirido (se considera común, por mitades) o que haya \textbf{condominios}, porque se compraron bienes conjuntamente durante la separación, como sucede habitualmente. Si hay sociedad conyugal, se trata de partes indivisas gananciales; si hay separación de bienes, de un condominio común.

El Código prevé la división del condominio durante el matrimonio, salvo que haya un interés familiar comprometido (por ejemplo, si el condominio es el hogar conyugal con hijos). La liquidación y partición de herencias en este régimen se limita a esos casos: condominios entre cónyuges o bienes cuya titularidad no puede determinarse. Se hará en el ámbito de la sucesión, es decir, en la partición y no mediante la división del condominio, dentro del régimen de la comunidad hereditaria; primero habrán de satisfacerse las deudas. A la expositora la remisión del régimen de separación de bienes a la partición no le parece acertada: considera que deberían ir por la vía de la división del condominio.

En consecuencia, en la separación de bienes \textbf{el cónyuge hereda como un hijo más} y no hay exclusión, porque no hay bienes gananciales. Respecto del antiguo art. 3576, la regla sigue igual: si la sociedad conyugal se disuelve por muerte y hay solo bienes gananciales, el cónyuge queda excluido de ellos, que se reparten entre los hijos.

\section{Sociedad conyugal disuelta en vida: tres hipótesis}

La doctrina actual distingue tres situaciones de sociedad conyugal disuelta en vida. El CCyC \textbf{no asumió ninguna de las posiciones doctrinarias}, de modo que seguirán discutiendo estos temas porque el código guardó silencio.

\begin{enumerate}
    \item \textbf{Disuelta y liquidada en vida.} Puede darse ahora al pasarse a un régimen convencional de separación de bienes. Ya no hay bienes gananciales: los que pertenecían al causante son bienes propios, y el cónyuge hereda como un hijo más. Una parte de la doctrina cuestiona que un cónyuge separado de bienes tenga preferencia frente a un cónyuge cuya sociedad se disuelve por la muerte; existen las dos posiciones. Desde el punto de vista legal, si la sociedad fue disuelta y liquidada en vida, no hay bienes gananciales sobre los cuales excluir al cónyuge.
    \item \textbf{Disuelta en vida pero no liquidada.} Se liquida junto con la división hereditaria; existe una \textbf{indivisión poscomunitaria}. Hay que diferenciar bienes propios y gananciales, y el cónyuge queda excluido por los hijos porque los bienes no liquidados todavía son gananciales.
    \item \textbf{Bienes adquiridos después de la disolución.} Cuando se disolvió la sociedad conyugal en vida y rige un régimen de separación de bienes, los bienes que antes habrían sido gananciales ya no lo son: son bienes personales del cónyuge, y el cónyuge supérstite hereda como un hijo más.
\end{enumerate}

\begin{table}[H]
\centering
\small
\begin{tabularx}{\textwidth}{
    >{\raggedright\arraybackslash}p{4.4cm}
    >{\raggedright\arraybackslash}X
    >{\raggedright\arraybackslash}p{3.6cm}
}
\toprule
\textbf{Situación al fallecer el causante} & \textbf{Bienes en juego} & \textbf{Consecuencia para el cónyuge supérstite} \\
\midrule
Sociedad conyugal disuelta y liquidada en vida (p.\,ej., pasaje a la separación de bienes por convención) & Ya no hay bienes gananciales: son bienes propios (personales) del causante & Hereda como un hijo más (la doctrina discute la solución) \\
\addlinespace
Disuelta en vida pero no liquidada (indivisión poscomunitaria) & Se distinguen bienes propios y gananciales; se liquida junto con la partición hereditaria & Los hijos lo excluyen respecto de los gananciales (aún no liquidados) \\
\addlinespace
Bienes adquiridos después de disuelta la sociedad conyugal (régimen de separación) & Bienes personales del causante, no gananciales & Hereda como un hijo más \\
\bottomrule
\end{tabularx}
\caption*{Sociedad conyugal disuelta en vida y herencia del cónyuge.}
\end{table}

\section{Cónyuge putativo, bigamia y caducidad de la acción de nulidad}

\subsection{El cónyuge putativo}

El CCyC mantiene lo relativo a la nulidad de matrimonio y nada dice sobre el cónyuge putativo.

\begin{itemize}
    \item Si la nulidad se produce \textbf{antes} de la muerte, no hay título y, por lo tanto, no hay vocación hereditaria.
    \item Si la nulidad se produce \textbf{después} de la muerte, rige el matrimonio putativo: el \textbf{cónyuge de buena fe hereda}, porque el matrimonio es válido y no hay efecto retroactivo debido a la buena fe. Si no hay buena fe, el matrimonio no produce efecto alguno y no hay título para heredar.
    \item Si no hay sentencia de nulidad, existe un título oponible \textit{erga omnes} y, hasta que no se declare la nulidad, el cónyuge viene a la sucesión.
\end{itemize}

\subsection{Bigamia y concurrencia con el cónyuge legítimo}

Todos coinciden en que el cónyuge de buena fe, cuya nulidad se produce después de la muerte, hereda porque tiene título. Distinto es el caso de \textbf{bigamia}, en que el cónyuge de buena fe concurre con el cónyuge legítimo. Esto \textbf{no está resuelto por el CCyC} y la doctrina se divide:

\begin{itemize}
    \item una posición sostiene que el cónyuge putativo de buena fe hereda junto con el legítimo \textbf{por partes iguales} (concurrencia);
    \item otra posición ---a la que adhiere la expositora--- sostiene que el \textbf{legítimo excluye al putativo}. El putativo podría heredar cuando el legítimo no viene porque está excluido por otras causas o porque lo excluyen los herederos (por ejemplo, los hijos en los bienes gananciales), pero no cuando concurre el legítimo.
\end{itemize}

Esta posición se apoya en la diferenciación que hacía el art. 1316 (liquidaciones de sociedades conyugales sucesivas, que daba preeminencia al cónyuge legítimo). Hoy el \textbf{art. 504} del CCyC sigue dando preferencia al cónyuge legítimo, por lo que se repetirá la discusión.

Aun cuando no concurra con el legítimo, el cónyuge putativo tiene, por el \textbf{art. 429, inc. c)}, una \textbf{acción de daños y perjuicios}, la misma que otorgaba el art. 225 del Código Civil, que puede resultar mucho más amplia que venir a la herencia, según los bienes que haya en ella. Así, el cónyuge putativo siempre está cubierto: viene a la herencia según una posición doctrinaria o, si no viene, tiene la acción de daños y perjuicios.

\subsection{Caducidad de la acción de nulidad por muerte (art. 239 de Vélez; art. 714 CCyC)}

El art. 239, siempre muy criticado, establece la \textbf{caducidad de la acción de nulidad de matrimonio por muerte} de uno o de los dos cónyuges. Con Vélez era absoluta; luego distintas leyes introdujeron excepciones: en caso de bigamia para el primer cónyuge del bígamo y para el segundo cónyuge de buena fe, y para ascendientes y descendientes en casos de nulidades absolutas cuando deben defender un interés propio (esto último viene del derecho canónico y, para la expositora, no tiene sentido).

La doctrina siempre sostuvo que el art. 239 debía modificarse para permitir ejercer la acción de nulidad después de la muerte de los cónyuges, aun en los casos de nulidades relativas, cuando exista un interés propio. El nuevo código, en el \textbf{art. 714}, reprodujo el art. 239: caducidad como principio general en caso de muerte de los cónyuges, con las excepciones que establecieron las leyes 2393, 17.711 y 23.515, todas surgidas de injusticias en casos concretos; todavía puede haber injusticias en algunos casos.

\section{El cónyuge con los ascendientes e indignidad de los ascendientes}

No ha habido modificaciones en el modo en que concurre el cónyuge con los ascendientes: concurre en \textbf{partes iguales}, sean los bienes propios o gananciales del causante que entraron a la herencia. La mitad es para el cónyuge y la mitad se reparte entre los ascendientes, sea uno, dos o cuatro.

La ley 23.264, al igualar a todos los ascendientes, legítimos o ilegítimos, fue muy criticada (antes solo venía a la herencia el padre natural). El código actual sigue la misma tesitura, pero quedan excluidos los ascendientes declarados indignos. En el Código de Vélez esto estaba en el art. 3296 bis y en otro artículo (el 249) que los excluía de la herencia cuando no hubieran reconocido al hijo en la minoría. Hoy todo se engloba en el artículo de indignidad: el \textbf{art. 2281, incs. e), f) y g)} prevé los mismos casos, es decir, el no reconocimiento durante la minoridad, no haber dado alimentos o haber sido privado de la autoridad parental (la antigua patria potestad).

Están igualados todos los ascendientes, pero aquellos que no cumplieron sus deberes no heredan: se mantiene el mismo sistema de la ley 23.264.
% TODO: la nota original cita «3596 bis» y las notas de revisión indican «3296 bis» como corrección editorial; verificar en la fuente.

\section{Derechos reconocidos al cónyuge supérstite}

El cónyuge hereda con los descendientes por cabeza en los bienes propios y queda excluido en los gananciales; con los ascendientes hereda por mitades y excluye a los colaterales.

\begin{table}[H]
\centering
\small
\begin{tabularx}{\textwidth}{
    >{\raggedright\arraybackslash}p{2.6cm}
    >{\raggedright\arraybackslash}X
    >{\raggedright\arraybackslash}X
}
\toprule
\textbf{Concurre con\ldots} & \textbf{Bienes propios} & \textbf{Bienes gananciales} \\
\midrule
Descendientes & Hereda por cabeza, como un hijo más & Queda excluido por los descendientes (art. 2433) \\
\addlinespace
Ascendientes & Mitad para el cónyuge y mitad para los ascendientes, sean uno, dos o cuatro & Ídem: partes iguales, cualquiera sea el origen de los bienes \\
\addlinespace
Colaterales & Los excluye & Los excluye \\
\bottomrule
\end{tabularx}
\caption*{Cómo concurre el cónyuge.}
\end{table}

En el nuevo código el cónyuge es legitimario y tiene una legítima de un medio, que \textbf{no ha sido disminuida}: es la misma que en el Código Civil. Tiene la investidura de pleno derecho, como los ascendientes y descendientes, y una preferencia en la administración de la sucesión cuando no hay mayoría. Se le reconocen, además, los siguientes derechos.

\subsection{Derecho real de habitación viudal (art. 2383)}

Es un derecho sobre un inmueble de propiedad del causante: no se refiere a un solo inmueble, sino a un inmueble de propiedad del causante. Opera de pleno derecho, es inoponible a los acreedores, no procede si el inmueble está en condominio con otros y exige que haya constituido el hogar conyugal. Debe ser de propiedad del causante y no estar en condominio con otra persona.

\subsection{Oposición a la partición de la vivienda (art. 2332)}

El art. 2332 permite al cónyuge oponerse a la partición respecto de la vivienda en que residieron los cónyuges, adquirida o constituida total o parcialmente con bienes gananciales. Lo tiene \textbf{de por vida}, con dos excepciones: los otros herederos pueden oponerse si hay otros bienes que cubran las necesidades del cónyuge, y este puede pedir que la vivienda se le adjudique en su lote.

Ambos artículos deberán coordinarse. El art. 2383 habla de un inmueble de propiedad del causante; el art. 2332, de una vivienda adquirida total o parcialmente con bienes gananciales, lo cual no impide que sea de propiedad del causante (puede haberse adquirido con bienes gananciales de propiedad del causante). No se sabe exactamente si en el primer caso, como la norma no distingue, quedan comprendidos bienes propios y gananciales, mientras que en el segundo basta con la adquisición total o parcial con bienes gananciales. Como los dos supuestos se superponen, habrá que determinar si el cónyuge podrá pedir la adjudicación en su lote y, a la vez, el derecho de habitación viudal, o si ello será considerado un \textbf{abuso del derecho}.

\subsection{Atribución preferencial de bienes y prelegado}

Se otorga a los herederos y también al cónyuge el derecho a pedir la atribución de determinados bienes en su hijuela, pagando el saldo. El cónyuge podría pedir, por ejemplo, establecimientos industriales, comerciales, agrícolas o ganaderos que haya ayudado a formar, pero siempre deberá pagar el saldo.

Además, todo legado hecho a un legitimario ---en este caso, al cónyuge--- se considera un \textbf{prelegado}, es decir, una mejora, salvo que el causante haya dispuesto lo contrario. Así, se han reconocido al cónyuge mucho más derechos que actualmente para oponerse a la partición y para pedir de manera preferencial determinados bienes.

\section{Causales de exclusión conyugal: el matrimonio \textit{in extremis}}

Las causales de exclusión conyugal son un tema en el que sí hay modificaciones. En el Código Civil siempre hubo causales de exclusión que se fueron modificando conforme se modificaron la separación personal y el divorcio (al introducirse el art. 67 bis, las causales objetivas y el divorcio vincular). Actualmente se distinguen: el matrimonio \textit{in extremis}; las causales de separación personal y divorcio; y la separación de hecho sin voluntad de unirse, culpable.

Con el divorcio incausado, desaparecen todas las causales de separación personal, y el divorcio, unilateral o bilateral, siempre produce la exclusión. Tiene su paralelismo con la ley 23.515, que introdujo el divorcio.

\subsection{Matrimonio \textit{in extremis} (art. 3573 de Vélez; art. 2436 CCyC)}

El art. 3573 del Código Civil presenta una gran diferencia entre el artículo y la nota de Vélez: el artículo habla de la enfermedad y del fallecimiento por esa misma enfermedad dentro de los 30 días, sin referirse a su gravedad, mientras que la nota habla de matrimonio \textit{in extremis}, es decir, con un moribundo. Esto dio lugar a una gran discusión doctrinaria sobre la gravedad de la enfermedad. Además, la doctrina estableció que, junto con el elemento objetivo (la enfermedad), había un elemento subjetivo: el \textbf{conocimiento} de la enfermedad. Si alguien se casa con una persona que tiene una enfermedad mortal pero no lo sabe, no se da el supuesto de la norma ---la \textit{ratio legis}, que es la \textbf{captación de herencia}---, porque la exclusión responde a que el cónyuge se casó únicamente para captar la herencia.

El art. 2436 del CCyC incorporó estos dos requisitos, lo cual está bien: exige una enfermedad de tal gravedad que lleve directamente a la muerte dentro del plazo de 30 días ---un desenlace fatal previsible, ahora sí un matrimonio con un moribundo--- y que el cónyuge tenga conocimiento de esa situación.

El problema está en la \textbf{excepción}. El Código de Vélez original no traía ninguna, lo que dio lugar a grandes injusticias, pues podía haberse vivido en concubinato o en un noviazgo de muchos años, y lo que corresponde probar es que no se quiso captar la herencia. La ley 17.711 introdujo la excepción de \textbf{regularizar una situación de hecho}, que por su ambigüedad también generó problemas de interpretación: en principio concubinato, pero luego se admitió también un noviazgo, una situación afectiva, el reconocimiento de un hijo, es decir, situaciones que permitieran probar que no hubo captación de herencia. Las \textbf{Jornadas de Bariloche} consideraron que debía admitirse con toda amplitud la prueba de que no hubo captación, y que la expresión «regularizar una situación de hecho» era muy acotada.

El art. 2436 del CCyC, según la expositora, \textbf{empeora} la situación: exige que el matrimonio haya sido \textbf{precedido de una unión convivencial}, y esta requiere permanencia, exclusividad y \textbf{dos años}. Si hubo una unión de hecho de un año, el cónyuge queda excluido aunque pueda probar que no hubo captación de herencia. Siendo la norma una presunción \textit{iuris tantum} de captación de herencia (objetivamente, se prueba la enfermedad grave y la muerte dentro de 30 días), debería permitirse probar de distintas maneras que no la hubo.

\begin{table}[H]
\centering
\small
\begin{tabularx}{\textwidth}{
    >{\raggedright\arraybackslash}p{2.2cm}
    >{\raggedright\arraybackslash}X
    >{\raggedright\arraybackslash}X
}
\toprule
\textbf{Aspecto} & \textbf{Código de Vélez (art. 3573)} & \textbf{CCyC (art. 2436)} \\
\midrule
Enfermedad & El artículo habla solo de «enfermedad»; la nota, de matrimonio con un moribundo. La doctrina exigió gravedad. & Enfermedad de desenlace fatal previsible, que lleve directamente a la muerte dentro de 30 días. \\
\addlinespace
Conocimiento & No figuraba; la doctrina lo exigió (sin él no hay captación de herencia). & El cónyuge debe conocer la situación. \\
\addlinespace
Excepción & Ninguna en el texto original (grandes injusticias). La ley 17.711 agregó «regularizar una situación de hecho». & Que el matrimonio haya sido precedido de una unión convivencial (permanencia, exclusividad y dos años). \\
\addlinespace
Prueba contraria & Las Jornadas de Bariloche propusieron admitir toda prueba de que no hubo captación de herencia. & Solo la unión convivencial: una unión de hecho de un año no alcanza. La expositora considera que la norma «empeora». \\
\bottomrule
\end{tabularx}
\caption*{Comparación del matrimonio \textit{in extremis} en ambos códigos.}
\end{table}

\section{Exclusión por divorcio y por separación de hecho}

\subsection{El divorcio incausado y la ley aplicable en el tiempo}

Quedan excluidos todos los que estén divorciados, con sentencia firme de divorcio. Ya no hay distinción entre causales objetivas y subjetivas ni entre inocente y culpable: todos quedan excluidos con el divorcio incausado, unilateral o bilateral. Si un matrimonio puede durar una semana, se pregunta la expositora, \textquestiondown{}por qué no excluir con el divorcio?

El único tema que podría presentarse es el del \textbf{inocente en la separación personal}. El art. 8 de la ley 26.994 (que aprobó el CCyC) establece la posibilidad de que la separación personal se convierta en divorcio, pero nada dice sobre la separación personal no convertida. Algo similar ocurrió con la ley 23.515: la ley 14.394 permitió durante un año el divorcio y el art. 6 de la ley 17.711 permitió la herencia del inocente de ese divorcio vincular. Cuando entró en vigencia la exclusión por divorcio sin hacer distinción (ley 23.515), el problema llegó a los tribunales y hubo un \textbf{plenario de la Cámara Civil} (alrededor de los años 90) que consideró que el derecho subjetivo se adquiere con la muerte del causante.

En materia sucesoria, el problema de la temporalidad de la ley es menor, porque \textbf{la ley vigente al momento de la muerte rige la sucesión}. Por lo tanto, el inocente de la separación heredará si el causante muere antes de la entrada en vigencia del CCyC (1.\textordmasculine{} de agosto), pero no si muere después, porque no está contemplado en el Código Civil que se deroga. El derecho subjetivo se adquiere en el momento de la muerte. Existe un fallo de la Corte sobre la legitimación de los hijos extramatrimoniales y la aplicación de leyes anteriores, pero se basó en convenciones internacionales y no en la ley de la muerte; esas convenciones podrían servir ahora en los casos de la mejora de los incapaces y discapacitados. En definitiva, rige la ley del momento de la muerte: la transmisión se produce en ese instante y también se adquiere en ese momento el derecho subjetivo.

\subsection{Separación de hecho sin voluntad de unirse}

La causal de separación de hecho, vigente en el art. 3575 del Código Civil, dio lugar a una discusión doctrinaria de más de cien años. En el Código Civil argentino la separación de hecho \textbf{no tiene relevancia jurídica salvo que sea sin voluntad de unirse}: si es con voluntad de unirse (por una enfermedad o un viaje por estudio, por ejemplo), no tiene relevancia.

\begin{itemize}
    \item Una posición sostuvo que, con la separación de hecho, no hay más herencia para el cónyuge separado.
    \item Otra posición (Machado) sostuvo que lo relevante es la falta de voluntad de unirse. Ejemplo planteado: un hombre deja a su mujer y se va con otra; se le muere la otra y quiere volver con su mujer; si la mujer no lo acepta, el hombre se pone en voluntad de unirse y, si ella muere, heredaría. Para excluirlo, la mujer tendría que haber iniciado juicio de divorcio cuando él se fue; es decir, se exigía un juicio de divorcio para excluir al otro cónyuge.
    \item Otra parte de la doctrina sostuvo que el problema estaba en la \textbf{culpabilidad}: quien abandona al otro, \textquestiondown{}por qué habría de perder la vocación hereditaria el abandonado? Esto aparece con la ley 17.711. Aun así, una posición doctrinaria (la de Fuentes) consideraba que la culpabilidad estaba en la falta de voluntad de unirse y que lo que se evitaba eran los divorcios \textit{post mortem}, mientras que otra parte, al igual que en el divorcio, entendía que el culpable de la separación era quien no heredaba. Luego vino el tema de la carga de la prueba: en el plenario \textbf{«Mauri»}, quienes querían excluir debían probar la separación sin voluntad de unirse más la culpabilidad.
    \item Con la ley 23.515, el art. 204 (de mala técnica legislativa) estableció en su segunda parte que el inocente podía reservar sus derechos probando su inocencia ---no se puede hacer probar a nadie su inocencia tras 25 o 30 años de matrimonio, sino que se prueba la culpabilidad del otro---. Reaparece así, a través de la doctrina y de un fallo de la Sala C y de la Suprema Corte de Buenos Aires, la \textbf{inversión de la carga de la prueba}: en paralelo con el art. 204, se interpretó que la segunda parte del plenario «Mauri» está derogada, y quien quiera heredar, el cónyuge supérstite separado de hecho, debe probar su inocencia. Así está la situación actual.
\end{itemize}

En la jurisprudencia hay una confusión tremenda: desde aplicar a rajatabla el plenario «Mauri» hasta establecer una causal netamente objetiva o la inversión de la carga de la prueba; hay cuatro o cinco posiciones.

\begin{table}[H]
\centering
\small
\begin{tabularx}{\textwidth}{
    >{\raggedright\arraybackslash}p{3cm}
    >{\raggedright\arraybackslash}X
    >{\raggedright\arraybackslash}X
}
\toprule
\textbf{Etapa} & \textbf{Criterio} & \textbf{Observaciones} \\
\midrule
Código de Vélez (original) & Excluye la separación de hecho sin voluntad de unirse & Discusión: \textquestiondown{}basta la separación de hecho o lo relevante es la falta de voluntad de unirse? (Machado y otros) \\
\addlinespace
Ley 17.711 (art. 3575) & Se incorpora la culpabilidad: hereda quien no causó la separación & Plenario «Mauri»: quien quiere excluir prueba separación sin voluntad de unirse más culpabilidad \\
\addlinespace
Ley 23.515 (art. 204) & Separación de hecho sin voluntad de unirse como causal de separación personal y divorcio & Se invierte la carga de la prueba: el cónyuge supérstite debe probar su inocencia; la jurisprudencia es confusa (cuatro o cinco posiciones) \\
\addlinespace
Código Civil y Comercial & Se vacía de toda culpa: excluye la separación de hecho sin voluntad de unirse & Sin inocentes ni culpables y sin plazos: ambos quedan excluidos; el abandonado tampoco puede pedir el divorcio para excluir al otro \\
\bottomrule
\end{tabularx}
\caption*{Evolución de la separación de hecho como causal de exclusión.}
\end{table}

\subsection{Lo que trae el Código Civil y Comercial}

En las relaciones matrimoniales el CCyC está \textbf{vaciado de toda culpa}, lo que, para la expositora, va en detrimento del inocente. Se vuelve al criterio del código original: la separación de hecho sin voluntad de unirse excluye, el abandonado queda excluido y no se han fijado plazos. La expositora sostiene que habrá una verdadera desheredación: quien se va deshereda, sin culpa, al que se queda. Basta la separación de hecho, pero la interpretación de la falta de voluntad de unirse dará la nota de \textbf{permanencia}: si se fue hace una semana, evidentemente la separación excluyente de la vocación hereditaria conyugal debe ser más o menos permanente. El código no lo dice ni establece plazo alguno, por lo que dará lugar a distintas posiciones doctrinarias, como en el código original.

Antes el abandonado podía, según la posición extrema, pedir un divorcio y excluir al otro; ahora no podrá pedir ningún divorcio porque quedará excluido igual. En definitiva, \textbf{ambos quedan excluidos en caso de separación de hecho}: no hay más inocentes ni culpables.

\begin{important}
\textbf{Exclusión del cónyuge en el nuevo Código.}
\begin{itemize}
    \item \textbf{Divorcio} (unilateral o bilateral): excluye siempre; ya no hay inocentes ni culpables.
    \item \textbf{Separación de hecho sin voluntad de unirse:} excluye a ambos cónyuges, sin plazo.
    \item \textbf{Matrimonio \textit{in extremis}:} solo se salva quien pruebe una unión convivencial previa (con permanencia, exclusividad y dos años).
\end{itemize}
\end{important}

\subsection{Reflexiones finales de la expositora sobre el matrimonio}

La expositora considera que se ha disminuido todo lo que es el matrimonio: se le ha sacado el concepto de \textbf{permanencia}. Observa que los primeros en aceptar el divorcio incausado fueron el código sueco y los del norte de Europa, y que todos establecen plazos para conseguirlo, hasta de un año de separación. Nuestro código ni siquiera exige los tres meses del código español, que fue su fuente. A su juicio, es más fácil disolver un matrimonio que resolver un contrato; y las uniones convivenciales tienen más derechos y deberes que el matrimonio. La parte de legítima será dada por otro docente.

\section{Cuadro Cornell: sucesión del cónyuge}

\begin{longtable}{
    >{\raggedright\arraybackslash}p{0.27\textwidth}
    >{\raggedright\arraybackslash}p{0.67\textwidth}
}
\toprule
\textbf{Preguntas / palabras clave} &
\textbf{Notas / respuestas} \\
\midrule
\endfirsthead
\toprule
\textbf{Preguntas / palabras clave} &
\textbf{Notas / respuestas} \\
\midrule
\endhead
\bottomrule
\endfoot

\textbf{\textquestiondown{}Cómo hereda el cónyuge en separación de bienes?} &
El Código no lo regula expresamente: por analogía hereda como en los bienes propios, es decir, como un hijo más. Solo puede haber problemas de determinación de la propiedad o condominios, que se resuelven en la partición (art. 508). \\[0.6em]

\textbf{\textquestiondown{}Qué dice el art. 2433?} &
Reúne los antiguos arts. 3570 y 3576: sobre los bienes propios el cónyuge hereda por cabeza con los hijos; sobre los gananciales lo excluyen los descendientes. \\[0.6em]

\textbf{\textquestiondown{}Qué pasa con la sociedad conyugal disuelta en vida?} &
Si está liquidada, hereda como un hijo más (bienes propios). Si no está liquidada, se distinguen propios y gananciales y los hijos lo excluyen de los gananciales. Los bienes adquiridos después de la disolución son personales. El Código guarda silencio y subsisten las posiciones doctrinarias. \\[0.6em]

\textbf{\textquestiondown{}Hereda el cónyuge putativo?} &
Sí, si es de buena fe y la nulidad se declara después de la muerte. En caso de bigamia, la doctrina se divide entre concurrencia por partes iguales o exclusión por el cónyuge legítimo (art. 504); además tiene acción de daños y perjuicios (art. 429). \\[0.6em]

\textbf{\textquestiondown{}Qué derechos especiales tiene el cónyuge supérstite?} &
Legítima de un medio, investidura de pleno derecho, preferencia en la administración, habitación viudal (art. 2383), oposición a la partición de la vivienda (art. 2332), atribución preferencial de bienes pagando el saldo y prelegado (todo legado al cónyuge se considera mejora). \\[0.6em]

\textbf{\textquestiondown{}Cuándo hay matrimonio \textit{in extremis}?} &
Cuando una enfermedad de desenlace fatal previsible lleva a la muerte en 30 días y el cónyuge lo sabía (art. 2436). La única excepción es la unión convivencial previa, lo que la expositora considera un empeoramiento. \\[0.6em]

\textbf{\textquestiondown{}Quién queda excluido por divorcio o separación de hecho?} &
Con el divorcio incausado, ambos cónyuges quedan excluidos. La separación de hecho sin voluntad de unirse también excluye a ambos, sin culpables ni plazos; se pierde la protección al inocente o abandonado. \\[0.6em]

\midrule
\multicolumn{2}{p{\dimexpr0.94\textwidth+2\tabcolsep\relax}}{
\textbf{Resumen:} Las modificaciones más importantes se concentran en el cónyuge. El divorcio incausado vacía de culpa las causales de exclusión, se endurece el matrimonio \textit{in extremis} y la separación de hecho excluye a ambos; a la vez, se le reconocen más derechos (habitación viudal, oposición a la partición, atribución preferencial). Muchas cuestiones (separación de bienes, sociedad conyugal disuelta, bigamia) quedan libradas a la doctrina y a la jurisprudencia.
} \\
\end{longtable}

% ------------------------------------------------------------------
\chapter[La legítima]{La legítima: reformas y cálculo}

Este capítulo está a cargo del docente de la segunda parte (legítima). Por ser un tema extenso, y porque en la exposición anterior la parte de colación debió sintetizarse mucho, se comienza resumiendo las principales reformas del CCyC en materia de legítima para luego desarrollar cada una.

\section{Principales reformas: aciertos y desaciertos}

\subsection{El principal acierto: la reducción de las cuotas (art. 2445)}

El principal acierto del CCyC es la \textbf{reducción de las cuotas de legítima} del art. 2445, una propuesta que venía haciendo la doctrina para ampliar la parte disponible y la libertad de testar. La legítima de los descendientes pasó de $4/5$ a $2/3$, con la posibilidad de reducirse aún más si hay un heredero con discapacidad (reduciendo el tercio de mejora, la legítima puede quedar en un tercio). La de los ascendientes pasó de $2/3$ en el Código de Vélez a $1/2$. Se mantuvo la cuota del cónyuge en $1/2$.

\begin{table}[H]
\centering
\small
\begin{tabularx}{\textwidth}{
    >{\raggedright\arraybackslash}p{3cm}
    >{\raggedright\arraybackslash}p{4.3cm}
    >{\raggedright\arraybackslash}X
}
\toprule
\textbf{Legitimario} & \textbf{Código de Vélez} & \textbf{CCyC (art. 2445)} \\
\midrule
Descendientes & $4/5$ & $2/3$ (puede reducirse a $1/3$ si hay un heredero con discapacidad: art. 2448) \\
\addlinespace
Ascendientes & $2/3$ & $1/2$ \\
\addlinespace
Cónyuge & $1/2$ & $1/2$ (se mantiene) \\
\addlinespace
Nuera viuda & Derecho sucesorio; legítima de $1/5$ & Derecho derogado: queda sin porción legítima \\
\bottomrule
\end{tabularx}
\caption*{Cuotas de legítima.}
\end{table}

Se ha derogado el derecho sucesorio de la nuera viuda (sin hijos o cuyos hijos no vivieran al abrirse la sucesión de los suegros). Su legítima era la cuarta parte de los bienes que hubiera recibido el marido si viviera; como el marido tenía $4/5$ de legítima, la nuera viuda tenía un quinto. Al derogarse su derecho sucesorio, también quedó sin porción legítima.

\subsection{Reformas opinables}

\paragraph{Mejora al heredero con discapacidad (art. 2448).} Es un fin loable, pero según las circunstancias de cada sucesión puede significar cierta desprotección de los restantes legitimarios. Por cualquier medio, incluso mediante fideicomiso, se puede disponer, además de la porción disponible, de otro tercio de la porción legítima para aplicarlo como mejora estricta a descendientes o ascendientes con discapacidad. Se sigue de algún modo el sistema español, que divide en tercios: uno de libre disposición, uno de mejora para alguno de los legitimarios y uno de legítima. El art. 2448 \textbf{omite al cónyuge}: el cónyuge con discapacidad no podrá tener esta mejora, aunque tendrá otras posibilidades según el régimen patrimonial matrimonial.

\paragraph{Donación inoficiosa al legitimario: colación y reducción (art. 2386).} Se recepta la tesis de Zannoni (expresada en la revista \textit{Derecho Privado Comunitario}), que la expositora considera adecuada. La donación inoficiosa al legitimario, es decir, la que excede la parte de libre disposición, \textbf{no solo está sujeta a colación sino también a reducción}. Esto estaba discutido y ahora es disposición expresa del art. 2386, ubicado en materia de colación pero referido a la protección de la legítima. La reducción alcanzará así no solo a las donaciones a terceros sino también a las efectuadas a legitimarios cuando excedan la parte disponible.

\subsection{Tres reformas desafortunadas}

\paragraph{Primera: el cálculo con valores heterogéneos.} Según el art. 2445, para el cálculo de la legítima se toma el valor líquido de los bienes hereditarios sumadas las donaciones según el valor que tenían a la época de la donación llevada a la época de la partición; esto equivale a sumar manzanas con peras y es contrario a la noción básica de legítima: \textbf{la legítima se calcula a la muerte, no a la partición}. Se confunde el cálculo de la masa de legítima con el de la masa partible. Belluscio ya había planteado que el donatario (por ejemplo, un tercero extraño) no tenía por qué esperar al cese de la indivisión y a la partición para conocer si su donación es inoficiosa y pasible de reducción. La legítima siempre se calculó al momento de la apertura de la sucesión. Las 24.\textordfeminine{} Jornadas Nacionales de Derecho Civil, celebradas en la misma facultad y presididas por el expositor junto con la doctora Hernández y el doctor Natale, plantearon por unanimidad que había un error en el cálculo, lo cual fue desoído.

\paragraph{Segunda: la desprotección del legitimario.} Azpiri, partidario de la libre disposición y no de la legítima, lo señala en su libro de \textit{incidencias}: si se mantiene la legítima, hay que protegerla. En materia de reducción y de efectos frente a terceros subadquirentes, el art. 2459 incorpora una \textbf{prescripción adquisitiva breve} de posesión por diez años de la cosa donada, que puede enervar la acción de reducción aun antes de que nazca: las acciones sucesorias nacen a partir de la muerte, mientras que la posesión de diez años puede cumplirse en vida del causante. Como dijo Zannoni, si un día antes de la muerte el donante escritura una donación de un inmueble que se poseía desde hace diez años, la acción no es viable. Algunos escribanos dirán que es la forma de sanear títulos y agilizar el tráfico jurídico, y es verdad que se tuvo esa mira, pero se limita seriamente la posibilidad de iniciar la reducción e incluso de calcular la legítima, pues ya no habrá que computar las donaciones efectuadas más de diez años atrás.

\paragraph{Tercera: la contradicción entre los arts. 2461 y 2449.} El art. 2461, en su parte final, y el art. 2449 generan una contradicción evidente. El art. 2461 se refiere a la transmisión de bienes a legitimarios (hoy art. 3604 de Vélez): cuando el causante transmite a un legitimario por contrato la propiedad de un bien con reserva de usufructo, la ley presume una gratuidad dentro de ese acto oneroso y también una mejora; se considera que hay una donación con imputación a la parte disponible y colación del resto (lo que los franceses llamaban colación para la reducción, y no para la igualdad). Las contradicciones se detallan en el capítulo 5 (usufructo, renta vitalicia y transmisión de bienes a legitimarios); en resumen: se afirma una presunción \textit{iure et de iure} de gratuidad pero se permite deducir lo efectivamente pagado, y se admite que los demás legitimarios consientan la enajenación, con lo que renuncian a la colación y a la reducción, es decir, a la legítima futura, en abierta contradicción con el art. 2449 (\textquestiondown{}se puede renunciar a la legítima futura o no? \textquestiondown{}es un pacto sobre herencia futura admitido o no?).

\section{Naturaleza de la legítima: \textit{pars bonorum} o \textit{pars hereditatis}}

El CCyC podía haber solucionado en la noción de legítima las discusiones sobre las teorías de la \textit{pars bonorum} y la \textit{pars hereditatis}; no lo hizo. Los partidarios de una y otra seguirán argumentando con nuevos elementos. En la definición del art. 2444 se reitera el artículo del proyecto de 1998 redactado por la doctora María José Méndez Costa, seguidora de la \textit{pars bonorum}, pero se suprime la referencia a los bienes.

\begin{formula}{Art. 2444 CCyC --- Legitimarios}
«Tienen una porción legítima de la que no pueden ser privados por testamento ni por actos de disposición entre vivos a título gratuito, los descendientes, los ascendientes y el cónyuge.»
\end{formula}

Algunos sostienen que, al suprimirse la parte del proyecto de 1998 («una porción legítima de los bienes»), se reafirma el carácter de parte de la herencia: no se puede recibir la legítima sin aceptar la herencia ni sin revestir la condición de legitimario. Pero, además, en todo el título no se habla de herederos forzosos: si no se habla de ellos, es porque \textbf{no es necesario ser heredero para ser legitimario}. En la preterición del legitimario (testamento con omisión de algún legitimario) se dice que la acción de entrega de legítima al preterido es a título de heredero de cuota, con lo cual no tiene derecho a acrecer, no tiene condición hereditaria y es un legitimario no heredero: se está de nuevo en la \textit{pars bonorum}, con una acción complemento prácticamente similar al art. 3600.

\begin{formula}{Art. 2451 CCyC --- Complemento}
«El legitimario a quien el testador le ha dejado por cualquier título menos que la porción legítima, solo puede pedir el complemento.»
\end{formula}

Según Belluscio, que sigue la tesis de la \textit{pars bonorum}, «por cualquier título» puede significar un título no hereditario: un legitimario donatario o acreedor de cargos, entre otras posibilidades. En resumen, \textbf{la discusión seguirá existiendo}.

\begin{table}[H]
\centering
\small
\begin{tabularx}{\textwidth}{
    >{\raggedright\arraybackslash}p{3cm}
    >{\raggedright\arraybackslash}X
    >{\raggedright\arraybackslash}X
}
\toprule
 & \textbf{\textit{Pars hereditatis}} & \textbf{\textit{Pars bonorum}} \\
\midrule
Idea central & La legítima es una parte de la herencia & La legítima es una parte de los bienes \\
\addlinespace
Argumentos en el art. 2444 y siguientes & Se suprimió la referencia a «los bienes»; no se puede recibir la legítima sin aceptar la herencia & El título no habla de «herederos forzosos»: no hace falta ser heredero para ser legitimario \\
\addlinespace
Acciones & --- & El preterido pide la entrega de legítima «a título de heredero de cuota», pero el legitimario puede ser solo donatario o acreedor de cargos (art. 2451) \\
\addlinespace
Resultado & La discusión seguirá existiendo: el Código no la resolvió & La discusión seguirá existiendo: el Código no la resolvió \\
\bottomrule
\end{tabularx}
\caption*{\textit{Pars hereditatis} y \textit{pars bonorum} en el nuevo Código.}
\end{table}

\subsection[Otras observaciones]{Otras observaciones: discapacidad, nuera viuda y desheredación}

Además de admitirse la mejora del heredero con discapacidad para el descendiente o el ascendiente, pero no para el cónyuge, lo cual no se explica, tampoco se siguió el criterio de la ley 26.657 de protección de la salud mental ni de la ley 26.378, que ratificó la Convención de los Derechos de las Personas con Discapacidad; podrá haberse unificado el texto para no dejar tanta amplitud judicial en la determinación de los casos en que procede esta mejora del tercio. Ya se vio que se deroga el derecho de la nuera viuda y que, como se explicó con otro docente de la cátedra, tampoco hay desheredación: la herramienta que tenía el heredero forzoso para sancionar a quien lo ofendió no se mantuvo.

\section{Cálculo de la legítima hereditaria}

Según el art. 2445, la legítima se calcula sobre la suma del valor líquido de la herencia al tiempo de la muerte del causante, más el de los bienes donados computables para cada legitimario a la época de la partición, según el estado del bien a la época de la donación.

\begin{formula}{Art. 2445 CCyC --- Porciones legítimas (texto oficial)}
«La porción legítima de los descendientes es de dos tercios, la de los ascendientes de un medio y la del cónyuge de un medio.

Dichas porciones se calculan sobre la suma del valor líquido de la herencia al tiempo de la muerte del causante más el de los bienes donados computables para cada legitimario, a la época de la partición según el estado del bien a la época de la donación.

Para el cómputo de la porción de cada descendiente sólo se toman en cuenta las donaciones colacionables o reducibles, efectuadas a partir de los trescientos días anteriores a su nacimiento o, en su caso, al nacimiento del ascendiente a quien representa, y para el del cónyuge, las hechas después del matrimonio.»
\end{formula}

Se toma, entonces, el valor líquido de la herencia a la muerte, pero también se agregan las donaciones (a legitimarios o a terceros) por su valor a la época de la partición, considerando el estado del bien a la época de la donación. Como dijo la doctora Hernández, se confunden operaciones de cálculo de la legítima con partición hereditaria: \textbf{la legítima es una masa ideal de cálculo, totalmente independiente de la partición}. En la colación sí era comprensible tomar un valor constante para evitar que las fluctuaciones intrínsecas y extrínsecas del bien donado incidieran sobre la igualdad buscada, porque allí se está en una operación de partición. Pero la legítima hay que calcularla a la muerte.

\subsection{Por qué el cálculo no puede hacerse al tiempo de la partición}

\begin{enumerate}
    \item \textbf{Masa partible y masa de cálculo difieren} respecto de las donaciones a terceros que no sean inoficiosas: no ingresan a la partición, pero sí al cálculo.

\begin{example}{masa de cálculo $\neq$ masa partible}
El causante deja 70 pesos y donó 30 a un amigo. Esa donación está dentro de la parte disponible, por lo que no es inoficiosa y la partición se hace sobre 70. En cambio, la masa de cálculo de la legítima se hace sobre 100: se le agrega el valor de la donación.
\end{example}

    \item \textbf{La divisibilidad de la colación} también incide: la colación es una obligación divisible que puede ser pedida por unos legitimarios y no por otros. Siempre se calculan dentro de la masa de legítima todos los valores donados, pero luego se forman distintas masas de partición según los legitimarios hayan o no pedido la colación.
    \item \textbf{Los legados particulares} integran la hijuela de bajas de la partición, pero no se toman en cuenta para la masa de cálculo de la legítima.

\begin{example}{el legado al diariero de la esquina}
El testador, con un patrimonio de 100, lega los 100 al diariero de la esquina. Ese legado excede la parte disponible si hay legitimarios. La legítima se calcula sobre 100 sin computar el legado: el legado integra la hijuela de bajas, pero no integra el cálculo de la legítima. No es como las donaciones, que sí hay que traer en valor, sean a legitimarios o a terceros.
\end{example}

    \item \textbf{Los honorarios de los abogados}: si el cálculo se llevara al momento de la partición, los abogados deberían esperarla para cobrar sus honorarios, lo cual está en contra del art. 2359, que dispone expresamente que no tienen que esperar a la partición.
\end{enumerate}

\begin{important}
Los valores para el cálculo de la legítima siempre deben tomarse al momento de la \textbf{muerte}, no de la partición. Mezclar valores de distintas épocas es un cálculo erróneo.
\end{important}

\subsection{Exclusiones de la masa de cálculo}

Además de las donaciones a terceros no inoficiosas (que no integran la masa partible), hay que tener en cuenta las siguientes exclusiones:

\begin{itemize}
    \item \textbf{Donaciones anteriores a los 300 días previos al nacimiento} del descendiente (o del ascendiente a quien representa), según el art. 2445. Es similar a lo que hoy indica el art. 1832, inc. 1.\textordmasculine{}, del Código Civil, que habilita a pedir la reducción a los descendientes nacidos después de la donación cuando haya otros que existían al tiempo de la donación. Como mejor técnica legislativa, el CCyC habla de no computar las donaciones anteriores a los 300 días previos al nacimiento. Si no hay ningún descendiente, se puede donar lo que se quiera y no se computará luego para la masa de cálculo.
    \item \textbf{Donaciones anteriores al matrimonio} respecto del cónyuge: es aplicación del plenario «Espota» (22 de agosto de 2002), que privó de legitimación activa al cónyuge supérstite para demandar la colación de donaciones previas a la celebración de su matrimonio.
    \item \textbf{Actos del art. 2461 (antes art. 3604)} cuando todos los legitimarios reconocen la onerosidad de la transmisión de bienes del causante a alguno de ellos con reserva de usufructo o cargo de renta vitalicia: la operación se considera sincera y no hay que computar ningún acto gratuito con imputación a la parte disponible.
    \item \textbf{Bienes donados que perecen sin culpa del donatario} (art. 2455): no los alcanza la reducción, porque habrían perecido igual en cabeza del donante, y por ende tampoco se computan para el cálculo de la legítima.
\end{itemize}
% TODO: contenido incompleto o ambiguo en las notas originales: el apunte cita el art. 2454 al tratar el perecimiento de lo donado en el cálculo (el cuadro y la sección de efectos de la reducción lo ubican en el art. 2455).

\begin{table}[H]
\centering
\small
\begin{tabularx}{\textwidth}{
    >{\raggedright\arraybackslash}X
    >{\raggedright\arraybackslash}X
}
\toprule
\textbf{Se computa} & \textbf{No se computa} \\
\midrule
El valor líquido de la herencia a la muerte del causante (se deducen las deudas, no las cargas) & Las donaciones a descendientes anteriores a los 300 días previos a su nacimiento (o al de su ascendiente representado) \\
\addlinespace
Las donaciones computables a legitimarios y a terceros & Las donaciones al cónyuge anteriores al matrimonio (plenario «Espota») \\
\addlinespace
Los beneficios colacionables o reducibles (discutido el caso del fideicomiso, «Bogelius») & Los actos del art. 2461 cuando todos los legitimarios reconocen la onerosidad del acto \\
\addlinespace
--- & Las donaciones cuyo bien perece sin culpa del donatario (art. 2455) \\
\bottomrule
\end{tabularx}
\caption*{Qué entra y qué no entra en la masa de cálculo de la legítima.}
\end{table}

\paragraph{Beneficios colacionables y el caso «Bogelius».} El CCyC no es del todo claro respecto de los beneficios colacionables enumerados en colación, que no se vuelven a marcar al tratar la porción legítima. Al comienzo, el artículo dice que la porción legítima no puede ser privada por testamento ni por acto de disposición entre vivos a título gratuito (comprendiendo las donaciones y otros beneficios colacionables, aunque no sean estrictamente donaciones). Pero para el cálculo solo se prevé el valor líquido de la herencia más el de los bienes donados, lo que genera una discordancia y dudas sobre la aplicación del caso «Bogelius» (Cámara Civil, con el voto de Zannoni): se trataba de un trust celebrado en Inglaterra; Zannoni consideró que, más allá de que no implicara donación (hay una transmisión a un fiduciario y ciertos legitimarios como beneficiarios), era un acto de beneficio o liberalidad que debía ser colacionable y reducible.

\subsection{El «valor líquido»: \textquestiondown{}se deducen solo las deudas o también las cargas?}

El CCyC no resolvió qué se entiende por valor líquido de los bienes de la herencia, a pesar de que también fue recomendación de las Jornadas de 2013. Para el expositor, \textbf{se deben deducir solo las deudas}, aunque hay opiniones importantes en contra (Natale y Azpiri). Si lo que se calcula es la legítima al momento de la muerte, las cargas nacen después de la muerte: los honorarios del abogado, del partidor, del administrador, del inventariador y del tasador son todas obligaciones posteriores a la muerte, es decir, cargas. Así no habrá que esperar el cese de la indivisión para conocer los honorarios. Deben deducirse, por lo tanto, solo las deudas hereditarias, como admite la mayor parte de la doctrina y la jurisprudencia.

\subsection{Concurrencia de legitimarios (art. 2446)}

El art. 2446 recoge el criterio doctrinario de que la legítima menor se toma de la mayor. Si concurren solo descendientes o solo ascendientes, la porción disponible se calcula según la respectiva legítima; si concurren cónyuge y descendientes, la porción disponible se calcula según la \textbf{legítima mayor}: la mayor absorbe a la menor. No trata el caso de cónyuge y ascendientes porque ambos tienen la misma legítima y la parte disponible siempre es un medio.

\section{Protección de la legítima: reglas generales}

La protección de la legítima es el tema central. El Código recepta criterios generales que también estaban en el Código de Vélez (arts. 3598 y 3599), hoy regulados en los arts. 2447 y 2449: no se puede imponer condición ni gravamen a la legítima, y no se puede renunciar a la legítima futura.

\begin{formula}{Art. 2447 CCyC --- Intangibilidad de la legítima}
«El testador no puede imponer gravamen ni condición alguna a la porción legítima. Si lo hace, se tiene por no escrita.»
\end{formula}

\begin{formula}{Art. 2449 CCyC --- Irrenunciabilidad}
«Es irrenunciable la porción legítima de una sucesión no abierta.»
\end{formula}

El CCyC siguió prácticamente el proyecto de 1998; este decía que no eran válidas, y ahora se dice que se tienen por no escritas, pero el criterio protector es el mismo. La irrenunciabilidad coincide con el \textbf{art. 1010}: la herencia futura no puede ser objeto de contratos ni tampoco los derechos hereditarios eventuales sobre objetos particulares, con una excepción para los pactos referidos a una explotación productiva o a participaciones societarias, con miras a conservar la unidad de gestión o prevenir conflictos: pueden contener disposiciones referidas a futuros derechos hereditarios y establecer compensaciones si no afectan la legítima. Con ello se admiten los \textbf{protocolos de empresas familiares}, aunque se refieran a una herencia no deferida. De lo contrario, no es posible renunciar; solo se podrá renunciar a la legítima una vez deferida la herencia y abierto el derecho de opción.

Toda la regulación de la protección de la legítima está orientada a la parte testamentaria, y debe completarse con la reducción de las donaciones inoficiosas, reguladas en el Código de Vélez en los arts. 1830 y 1832 y ahora en el art. 1565.

\begin{formula}{Art. 1565 CCyC --- Donación inoficiosa}
«Se considera inoficiosa la donación cuyo valor excede la parte disponible del patrimonio del causante.» A este respecto se aplican los preceptos del Código sobre protección de la legítima.
\end{formula}

Estas normas generales (prohibición de condicionamiento, de gravamen y de renuncia a la legítima futura) darán pie a que siga existiendo jurisprudencia sobre \textbf{fraude genérico a la legítima} o fraude a través de la constitución de sociedades, con casos emblemáticos (el expositor menciona los casos «Astesiano» y «Morro Berner» si queda tiempo) y distintos supuestos considerados lesivos de la legítima.

\subsection{Los puntos débiles de la protección: la «protección relativa»}

Azpiri, en su libro de \textit{incidencias}, habla de cuatro o cinco aspectos que hacen a la desprotección de los legitimarios. Para él, la legítima se ha mantenido, pero con una suerte de \textbf{protección relativa}:

\begin{table}[H]
\centering
\small
\begin{tabularx}{\textwidth}{
    >{\centering\arraybackslash}p{0.9cm}
    >{\raggedright\arraybackslash}X
    >{\raggedright\arraybackslash}X
}
\toprule
\textbf{N.\textordmasculine{}} & \textbf{Disposición} & \textbf{Por qué debilita la protección} \\
\midrule
1 & Art. 2444: alude a testamentos y actos entre vivos a título gratuito; el cálculo solo menciona donaciones & Duda sobre beneficios colacionables como el fideicomiso (caso «Bogelius») \\
\addlinespace
2 & Art. 2459: prescripción adquisitiva de 10 años & Sanea títulos, pero reduce la masa de legítima y extingue la acción antes de nacer \\
\addlinespace
3 & Arts. 2283 y 2284: la indignidad puede oponerse en todo tiempo & Sin plazo de caducidad frente a la reducción, la colación o la petición de herencia \\
\addlinespace
4 & Art. 2461: el legitimario puede consentir la enajenación con reserva de usufructo, uso, habitación o renta vitalicia & Equivale a renunciar a la colación y la reducción, contra los arts. 1010 y 2449 \\
\addlinespace
5 & Art. 2448: mejora al heredero con discapacidad & Disminuye la masa de legítima y omite al cónyuge \\
\bottomrule
\end{tabularx}
\caption*{Puntos de «protección relativa» de la legítima (según Azpiri).}
\end{table}

En cuanto a la indignidad, el legitimario demandado por reducción, colación o petición de herencia puede oponer la indignidad sucesoria, para lo cual no hay plazo de caducidad (arts. 2283 y 2284); si se puede oponer en todo tiempo, hay menor protección. Respecto del art. 2459, Azpiri dice que sanea títulos pero disminuye la integración de la masa de legítima en perjuicio de los legitimarios. Respecto del art. 2448, aunque el fin es sumamente razonable, supone una disminución de la masa de legítima; y, como se dijo, se ha omitido al cónyuge.

\section{Cuadro Cornell: la legítima}

\begin{longtable}{
    >{\raggedright\arraybackslash}p{0.27\textwidth}
    >{\raggedright\arraybackslash}p{0.67\textwidth}
}
\toprule
\textbf{Preguntas / palabras clave} &
\textbf{Notas / respuestas} \\
\midrule
\endfirsthead
\toprule
\textbf{Preguntas / palabras clave} &
\textbf{Notas / respuestas} \\
\midrule
\endhead
\bottomrule
\endfoot

\textbf{\textquestiondown{}Qué cuotas de legítima fija el art. 2445?} &
Descendientes: $2/3$ (antes $4/5$); ascendientes: $1/2$ (antes $2/3$); cónyuge: $1/2$ (sin cambios). La nuera viuda pierde su derecho sucesorio y su legítima. \\[0.6em]

\textbf{\textquestiondown{}Qué es la mejora al heredero con discapacidad (art. 2448)?} &
Un tercio adicional de la legítima de descendientes o ascendientes con discapacidad (la legítima puede quedar en $1/3$). Omite al cónyuge con discapacidad. \\[0.6em]

\textbf{\textquestiondown{}Es el legitimario siempre heredero?} &
No hay respuesta definitiva en el Código: se suprimió la referencia a «los bienes» en el art. 2444 y no se habla de «herederos forzosos», por lo que la discusión entre \textit{pars hereditatis} y \textit{pars bonorum} seguirá existiendo. \\[0.6em]

\textbf{\textquestiondown{}Cuándo se calcula la legítima?} &
Al tiempo de la muerte, no de la partición. El art. 2445 confunde masa de cálculo con masa partible al valuar las donaciones a la época de la partición; el expositor lo considera un error. \\[0.6em]

\textbf{\textquestiondown{}Qué se excluye de la masa de cálculo?} &
Donaciones anteriores a los 300 días previos al nacimiento del descendiente, donaciones anteriores al matrimonio (plenario «Espota»), actos del art. 2461 con consentimiento de todos los legitimarios y bienes perecidos sin culpa del donatario (art. 2455). Se deducen las deudas, no las cargas. \\[0.6em]

\textbf{\textquestiondown{}Qué reglas protegen la legítima?} &
No admite gravamen ni condición (art. 2447) y es irrenunciable antes de la apertura (art. 2449; art. 1010 sobre herencia futura, con excepción de los protocolos de empresas familiares). \\[0.6em]

\midrule
\multicolumn{2}{p{\dimexpr0.94\textwidth+2\tabcolsep\relax}}{
\textbf{Resumen:} El acierto principal en materia de legítima es la reducción de las cuotas ($2/3$, $1/2$, $1/2$) para ampliar la libertad de testar. En cambio, el cálculo con valores heterogéneos, la prescripción adquisitiva de diez años del art. 2459 y las contradicciones del art. 2461 debilitan la protección del legitimario. La legítima se calcula a la muerte; es intangible e irrenunciable antes de la apertura de la sucesión.
} \\
\end{longtable}

% ------------------------------------------------------------------
\chapter[Acción de reducción]{Acción de reducción y sus efectos}

\section{La acción de reducción: concepto, orden y legitimación}

\begin{definition}{Acción de reducción}
Es la acción que tienen los legitimarios para atacar, en primer término, a los legados inoficiosos o, en el derecho actual, a los legados de cuota (a partir de la entrada en vigencia del CCyC, al heredero de cuota) y, en segundo término, a los donatarios o subadquirentes de bienes de estos donatarios que han recibido donaciones inoficiosas, en cuanto los beneficios recibidos excedan la parte disponible.
\end{definition}

Es la más importante de las acciones tutelares. Primero se reducen los legados y, a continuación, las donaciones.

\subsection{Reducción de disposiciones testamentarias (arts. 2450 y 2452)}

El art. 2452 habla de reducción de \textbf{disposiciones testamentarias}, lo que parece permitir que se reduzca también una institución de herederos. Sin embargo, no se puede reducir una institución de herederos, pues para ello existe una acción específica: la \textbf{acción de preterición del legitimario} o, como se la llama ahora, \textbf{acción de entrega de la legítima} (art. 2450), que es la misma acción del art. 3715 del Código de Vélez, algo peor regulada. Está empeorada porque no solo se da al legitimario omitido en el testamento, sino que su última parte se la otorga también al legitimario cuando el difunto no deja bienes pero efectuó donaciones, con lo que hay una evidente superposición con la acción de reducción: las donaciones siempre han sido atacadas por la vía de la reducción y no de la preterición.

A fin de recibir o complementar su porción, el legitimario afectado puede pedir la reducción de las instituciones de herederos de cuota y de los legados (art. 2452). Los legados se reducen «en el mismo orden» establecido en el segundo párrafo del art. 2358. El expositor considera que se \textbf{cortó y pegó} el mismo error que ya se había marcado en el proyecto de 1998: el orden de reducción de las disposiciones testamentarias (herederos de cuota y legados particulares) debe ser el \textbf{inverso} al orden de pago. De acuerdo con el art. 3795 del Código de Vélez, cuando la sucesión es insolvente se pagan primero los legados de cosa cierta, luego los remuneratorios de servicios y por último los de cantidad (de cuota).

En la primera parte de la norma se refleja el orden correcto: primero se reduce al heredero de cuota (que no tiene derecho a acrecer, salvo razones excepcionales). Pero luego dice que los legados se reducen en el mismo orden en que se pagan, es decir, el del art. 2358: primero, con la preferencia otorgada por el testador (muy razonable: el testador puede indicar qué legado reducir primero si se excede en la parte disponible); en segundo lugar, los de cosa cierta y determinada; luego los demás legados; si hay varios de la misma categoría, se pagan a prorrata. Reducir primero un legado de cosa cierta puede obligar a dividir o subdividir un inmueble, en lugar de reducir un legado remuneratorio de servicios: no tiene lógica, es antieconómico. Normalmente el orden de reducción ha sido el inverso al orden de los pagos; aquí se sigue el mismo.

\subsection{Reducción de las donaciones (art. 2453)}

Si la reducción de herederos de cuota y legados particulares no alcanza para cubrir la legítima, se reducen las donaciones inoficiosas. La reducción de donaciones solo opera si la de las disposiciones testamentarias no cubre la legítima, tesis triunfante en la jurisprudencia y la doctrina.

\begin{formula}{Art. 2453 CCyC --- Reducción de donaciones}
«Si la reducción de las disposiciones testamentarias no es suficiente para que quede cubierta la porción legítima, el heredero legitimario puede pedir la reducción de las donaciones hechas por el causante. Se reduce primero la última donación, luego las demás, en orden inverso a sus fechas, hasta salvar el derecho del reclamante. Las de igual fecha se reducen a prorrata.»
\end{formula}

Lo último es novedoso y le parece acertado al expositor: ha habido casos en que se celebraron varias escrituras de donación el mismo día, y no puede ocurrir que, por tener números correlativos, algunas sean reducidas y otras no, siendo que todas se efectúan en el mismo momento.

\begin{table}[H]
\centering
\small
\begin{tabularx}{\textwidth}{
    >{\raggedright\arraybackslash}p{3.2cm}
    >{\raggedright\arraybackslash}X
    >{\raggedright\arraybackslash}X
}
\toprule
\textbf{Concepto} & \textbf{Qué establece el Código} & \textbf{Observación del docente} \\
\midrule
Reducción de disposiciones testamentarias (art. 2452) & Primero el heredero de cuota; luego los legados, «en el mismo orden en que se pagan» (art. 2358) & Error copiado del proyecto de 1998: el orden de reducción debe ser el inverso al orden de pago \\
\addlinespace
Orden de pago (art. 2358) & 1.\textordmasculine{} con la preferencia otorgada por el testador; 2.\textordmasculine{} cosa cierta y determinada; 3.\textordmasculine{} los demás legados; los de igual categoría, a prorrata & Reducir primero un legado de cosa cierta puede obligar a dividir un inmueble: es antieconómico \\
\addlinespace
Reducción de donaciones (art. 2453) & Solo si no alcanza la reducción testamentaria; se reduce primero la última donación y luego las demás, en orden inverso a sus fechas & Las de igual fecha se reducen a prorrata: novedad que el docente considera acertada \\
\bottomrule
\end{tabularx}
\caption*{Orden de pago y orden de reducción.}
\end{table}

Por otro lado, el Código parece inclinarse por la acción real en el art. 2458, al hablar de reipersecución: el legitimario puede perseguir a terceros adquirentes de bienes registrables. Pero donatario y subadquirente pueden \textbf{desinteresar} al legitimario satisfaciendo en dinero el perjuicio de la cuota de legítima. Se otorga así la reipersecución propia de una acción real, pero se la relativiza con la posibilidad de desinteresar en dinero, posibilidad a la que ya se inclinaban la doctrina mayoritaria y alguna jurisprudencia. Las discrepancias subsisten; el Código no las solucionó.

\subsection{Legitimación activa y pasiva}

\paragraph{Legitimación activa.} Como novedad, la legitimación de los descendientes solo tiene lugar frente a donaciones que no superen los 300 días anteriores al nacimiento. Hay otra disminución de la legitimación activa frente a donaciones con un plazo de posesión superior a diez años: allí tampoco habrá legitimación. La primera limitación es lógica: si la persona no había sido concebida, no era posible el perjuicio por esa liberalidad. La segunda plantea el problema de la prescripción adquisitiva breve, que puede extinguir una acción sucesoria antes de nacer. También rige lo dicho sobre la legitimación del cónyuge y el plenario «Espota»: no pueden demandarse ni la colación ni la reducción de donaciones efectuadas antes del matrimonio.

\paragraph{Legitimación pasiva.} Depende del acto que se trate de atacar. Si se entabla contra un heredero de cuota o contra legatarios particulares, hay que tener en cuenta el orden de la reducción para ver si es viable; por ejemplo, no podría demandarse a un donatario antes de agotar la reducción de las disposiciones testamentarias. Respecto de los herederos de cuota es muy sencillo, pues será pasible de reducción la alícuota suficiente para dejar a salvo la legítima; los legados particulares también deben ser reducidos hasta dejar a salvo esa porción intangible.

Si hay que demandar a donatarios, ahora se puede demandar no solo a donatarios terceros sino también a \textbf{donatarios legitimarios}: según el art. 2386, el legitimario donatario que recibió una donación inoficiosa por exceder la parte de libre disposición está sujeto no solo a colación sino también a reducción, lo que al expositor le parece lógico y recoge la posición de Zannoni. El legitimario que demanda a otro legitimario puede quedar en peor posición, limitado por la acción personal de colación, que el legitimario que demanda a un tercero, a quien puede atacar con la vía reipersecutoria, aunque menguada, de la reducción.

\section{Efectos de la reducción (arts. 2454 a 2459)}

Los efectos de la acción de reducción tienen una regulación casuística y desglosada. Es preferible tenerla a no tenerla: en el Código de Vélez había prácticamente los arts. 3600 y 3601 y la norma del art. 3955, que habla indirectamente de la acción reipersecutoria o reivindicatoria dentro del contexto de una prescripción y que dio lugar a varias posiciones encontradas, que ahora se busca dejar de lado.

\begin{table}[H]
\centering
\small
\begin{tabularx}{\textwidth}{
    >{\raggedright\arraybackslash}p{5.5cm}
    >{\raggedright\arraybackslash}X
}
\toprule
\textbf{Jurisdicción} & \textbf{Interpretación del art. 3955 de Vélez} \\
\midrule
Capital Federal (Colegio de Escribanos, fallo «Escary c/ Pietranera», 1912) & Si la donación es a un heredero forzoso, se aplica la acción personal de colación; si es a un tercero, la reducción \\
\addlinespace
Provincia de Buenos Aires & Siempre acción personal (el art. 3955 queda casi como letra muerta) \\
\addlinespace
Entre Ríos & Siempre acción reivindicatoria, sean herederos forzosos o terceros los donatarios \\
\bottomrule
\end{tabularx}
\caption*{Las interpretaciones del art. 3955 del Código de Vélez.}
\end{table}

\subsection{Reducción total o parcial de la donación (art. 2454)}

Si la reducción es \textbf{total}, la donación queda \textbf{resuelta}: se aplica todo el valor de la donación a la inoficiosidad, por exceder la parte disponible y afectar la legítima (la condición resolutoria implícita de toda donación, sujeta a la inoficiosidad cuando median legitimarios). Si es \textbf{parcial}, por afectar solo en parte la legítima, hay dos posibilidades: si el bien es divisible, se divide entre el legitimario y el donatario; si es indivisible, la cosa queda para quien le corresponda una porción mayor, con un crédito a favor de la otra parte por el valor de su derecho.

Luego, el art. 2454 dispone que \textbf{en todo caso} el donatario puede impedir la resolución entregando al legitimario la suma de dinero necesaria para completar el valor de la porción legítima. El expositor señala que aquí aparecen contradicciones poco comprensibles, pues se venía hablando de resolución: es una resolución \textbf{menguada}, porque en lugar de entregar la cosa que, por exceder la parte disponible, afecta la legítima en su totalidad, puede entregarse dinero. El donatario es deudor de los \textbf{frutos} desde la notificación de la demanda (jurisprudencia pacífica); si opta por entregar dinero, debe los \textbf{intereses} desde el momento en que opta por enervar los efectos de la reducción, lo cual está bien porque el interés es un fruto del capital.

\subsection{Perecimiento de lo donado (art. 2455)}

El art. 2455 trae una solución novedosa: si el bien donado perece por \textbf{culpa} del donatario, este debe su valor. Si perece \textbf{sin culpa}, el valor de lo donado no se computa para el cálculo de la porción legítima y, por ende, no se reduce, pues habría perecido igual en cabeza del donante o del donatario. Si perece parcialmente por su culpa, debe la diferencia de valor; si perece parcialmente sin su culpa, se computa el valor subsistente. Habrá que atender a los casos en que, habiendo perecimiento, se cobra una indemnización sustitutiva, para que no queden desguarnecidos los legitimarios; de lo contrario, el expositor está de acuerdo con la solución.

\subsection{Derechos reales constituidos por el donatario (art. 2457)}

La reducción extingue, con relación al legitimario, los derechos reales constituidos por el donatario o por sus sucesores: si el donatario prendó el bien mueble, cae la prenda; si lo hipotecó, cae la hipoteca. Esto generará, a criterio del expositor, desprotección y discusiones, sobre todo con los beneficiarios de esos derechos reales.

\subsection{Acción reipersecutoria (art. 2458)}

\begin{formula}{Art. 2458 CCyC --- Acción reipersecutoria}
«El legitimario puede perseguir contra terceros adquirentes los bienes registrables. El donatario y el subadquirente demandados pueden desinteresar al legitimario satisfaciendo en dinero el perjuicio a la cuota legítima.»
\end{formula}

El art. 2458 parece limitarse a los bienes registrables. Sin embargo, en materia de muebles se recoge el principio del art. 2412 de Vélez: según el art. 1895, la posesión de buena fe de la cosa no hurtada ni perdida es suficiente para adquirir derechos reales, salvo que se pruebe la adquisición gratuita. Por lo tanto, también podrán perseguirse cosas muebles, en tanto se demuestre la mala fe o el título gratuito del adquirente. La reipersecución se ve menguada en la segunda parte de la norma: si el interés es la medida de la acción, en la medida en que el legitimario sea desinteresado con dinero, la reipersecución ya no procede.

\subsection{Prescripción adquisitiva breve de diez años (art. 2459)}

El art. 2459 es el más polémico y criticado por la doctrina hasta la fecha.

\begin{formula}{Art. 2459 CCyC --- Prescripción adquisitiva}
«La acción de reducción no procede contra el donatario ni contra el subadquirente que han poseído la cosa donada durante diez años computados desde la adquisición de la posesión.» Se aplica el art. 1901 (unión de posesiones: heredero del poseedor, cesionario, sucesor del poseedor).
\end{formula}

Se puede detener los efectos reipersecutorios de una acción de reducción alegando haber poseído la cosa donada por diez años. No se cuentan desde la muerte ni se limitan a los diez años anteriores al fallecimiento del causante (como decía el proyecto de 1998), sino \textbf{desde la adquisición de la posesión}: si el causante muere quince años después de donar, la acción de reducción está enervada por la prescripción adquisitiva del art. 2459. Es, como ya se dijo, una acción que se extinguió antes de nacer. El expositor considera que hay una evidente desprotección del legitimario, que se limita no solo el cálculo sino la procedencia misma de la reducción, y que generará problemas para el cálculo, pues se incluirá en la masa algo que luego no podrá reducirse.

\section{Usufructo, renta vitalicia y transmisión de bienes a legitimarios}

\subsection{Usufructo, uso, habitación y renta vitalicia (art. 2460)}

Lo que hoy regulan los arts. 3603 y 3604 del Código de Vélez (legados de usufructo y renta vitalicia; transmisión de bienes a legitimarios) estará regulado en los arts. 2460 y 2461.

\begin{formula}{Art. 2460 CCyC --- Constitución de usufructo, uso, habitación o renta vitalicia}
«Si la disposición gratuita entre vivos o el legado son de un usufructo, uso, habitación o renta vitalicia, el legitimario o, en su caso, todos los legitimarios de común acuerdo pueden optar entre cumplirlo o entregar al beneficiario la porción disponible.»
\end{formula}

La opción del art. 3603 para los legados de usufructo (cumplir el legado o entregar la parte disponible) está explicada por Vélez en la nota: el legado de usufructo no es reducible y por eso se da una opción que nunca perjudica al legitimario. El CCyC la mejora al agregar el uso y la habitación, y se agrega además la \textbf{donación} efectuada con reserva de usufructo, uso o habitación o con contraprestación de una renta vitalicia, hasta ahora no prevista expresamente (solo los legados).

Se suprime la mención del art. 3603 del legado de usufructo que excediera la parte de libre disposición: determinar cuándo un legado de usufructo excede la parte disponible depende de la vida del usufructuario. Vidal Taquini y Medina proponían acudir a las tablas de seguros de vida; la mayor parte de la doctrina decía que no había que hacer cálculos y que el legitimario era el juez de su interés. Esto también es un acierto.

Se mantiene la opción de cumplir o entregar la parte disponible. Se soluciona la exigencia de unanimidad (que con el art. 3613 se discutía: qué pasa si algunos legitimarios quieren cumplir el legado y otros entregar la parte disponible): ahora todos los legitimarios \textbf{de común acuerdo} pueden optar por una u otra alternativa. No dice qué ocurre si no hay unanimidad, por lo que subsisten las mismas discusiones: el expositor, con la doctora Hernández, sostiene que hay que cumplir el legado en caso de falta de unanimidad, siguiendo la tesis de Borda; otros dicen que hay que entregar la parte disponible; otros, que debe decidir el juez.

Tampoco menciona los legados y donaciones de \textbf{nuda propiedad}: solo las donaciones con reserva de usufructo y los legados con reserva de usufructo, uso, habitación y renta vitalicia. Si se considera que es una norma de excepción, la interpretación debería ser restrictiva, pero habrá que estar atentos a si la jurisprudencia la extiende.

\subsection{Transmisión de bienes a los legitimarios (art. 2461)}

La finalidad del art. 3604 era evitar que el causante, mediante una aparente o simulada compraventa, ocultase en realidad una donación cuando se reservaba el usufructo o imponía el cargo de una renta vitalicia. Hoy está regulado en el art. 2461, que en su primera parte ratifica este criterio.

\begin{formula}{Art. 2461 CCyC --- Transmisión de bienes a legitimarios}
«Si por un acto entre vivos a título oneroso el causante transmite a alguno de los legitimarios la propiedad de bienes con reserva de usufructo, uso o habitación, o con la contraprestación de una renta vitalicia, se presume sin admitir prueba en contrario la gratuidad del acto y la intención de mejorar al beneficiario.

Se debe deducir del valor de lo donado las sumas que el adquirente demuestre haber efectivamente pagado. El valor de los bienes debe ser imputado a la porción disponible y el excedente es objeto de colación. Esta imputación y colación no pueden ser demandadas por los legitimarios que consintieron en la enajenación, sea onerosa o gratuita, con alguna de las modalidades indicadas.»
\end{formula}

Se trata de una \textbf{presunción de gratuidad} (porque era un acto oneroso) y de una \textbf{presunción de mejora} (por la imputación primero a la parte disponible, que habilita la reducción del resto, que se colaciona). Es lo que la doctrina francesa distingue entre colación con fines de igualdad y colación con fines de reducción; esta sería una colación con fines de reducción.

La parte final del artículo trastoca el tema. Tras establecer una presunción \textit{iure et de iure}, se da la posibilidad de probar en contrario al deducir del valor de lo donado las sumas que el adquirente demuestre haber pagado: se abre una presunción \textit{iuris tantum}. Y lo peor es que lo que debía aplicarse solo a actos onerosos, por la última parte parece aplicarse también a los gratuitos: la imputación y la colación no pueden ser demandadas por los legitimarios que consintieron la enajenación, sea onerosa o gratuita, con alguna de las modalidades. Entonces, en una donación con reserva de usufructo en la que los restantes legitimarios presten consentimiento a la gratuidad del acto, estarán renunciando a pedir la colación y la reducción, pactando sobre una herencia futura y desprotegiéndose, en contradicción con el art. 2449.

\begin{table}[H]
\centering
\small
\begin{tabularx}{\textwidth}{
    >{\raggedright\arraybackslash}p{3.2cm}
    >{\raggedright\arraybackslash}p{3.6cm}
    >{\raggedright\arraybackslash}X
}
\toprule
\textbf{Aspecto} & \textbf{Art. 3604 (Vélez)} & \textbf{Art. 2461 (CCyC)} \\
\midrule
Modalidades & Reserva de usufructo o renta vitalicia & Se agregan el uso y la habitación \\
\addlinespace
Presunción & Gratuidad del acto y mejora (\textit{iure et de iure}) & Se ratifica, pero se deducen las sumas que el adquirente pruebe haber pagado: la presunción pasa a ser, en los hechos, \textit{iuris tantum} \\
\addlinespace
Consentimiento de los demás legitimarios & --- & No pueden demandar imputación ni colación quienes consintieron, sea el acto oneroso o gratuito: implica renunciar a la legítima futura, en contradicción con el art. 2449 \\
\bottomrule
\end{tabularx}
\caption*{Transmisión de bienes a legitimarios.}
\end{table}

\section{Cuadro Cornell: acción de reducción y sus efectos}

\begin{longtable}{
    >{\raggedright\arraybackslash}p{0.27\textwidth}
    >{\raggedright\arraybackslash}p{0.67\textwidth}
}
\toprule
\textbf{Preguntas / palabras clave} &
\textbf{Notas / respuestas} \\
\midrule
\endfirsthead
\toprule
\textbf{Preguntas / palabras clave} &
\textbf{Notas / respuestas} \\
\midrule
\endhead
\bottomrule
\endfoot

\textbf{\textquestiondown{}En qué orden se reduce?} &
Primero las disposiciones testamentarias (herederos de cuota y legados) y, si no alcanza, las donaciones, de la última a la primera (art. 2453), las de igual fecha a prorrata. El expositor critica que el orden de reducción de legados siga el orden de pago. \\[0.6em]

\textbf{\textquestiondown{}Cuáles son los efectos de la reducción de donaciones?} &
Reducción total: la donación queda resuelta; parcial: se divide el bien o se adjudica a quien tiene la porción mayor con crédito a la otra parte. El donatario puede enervarla pagando en dinero. Se extinguen los derechos reales del donatario (art. 2457) y hay acción reipersecutoria (art. 2458). \\[0.6em]

\textbf{\textquestiondown{}Por qué se critica el art. 2459?} &
Porque la posesión de diez años desde la adquisición enerva la acción de reducción antes de que nazca: protege títulos, pero desprotege al legitimario y limita el cálculo. \\[0.6em]

\textbf{\textquestiondown{}Qué contradicción presenta el art. 2461?} &
Presume \textit{iure et de iure} la gratuidad, pero admite deducir lo pagado y permite que los legitimarios consientan la enajenación, lo que implica renunciar a la legítima futura, en contra del art. 2449. \\[0.6em]

\midrule
\multicolumn{2}{p{\dimexpr0.94\textwidth+2\tabcolsep\relax}}{
\textbf{Resumen:} La reducción procede primero contra las disposiciones testamentarias y luego contra las donaciones, de la última a la primera. Sus efectos están regulados casuísticamente (resolución total o parcial, perecimiento, extinción de derechos reales, reipersecución) pero se ven menguados por la posibilidad de enervarla con dinero y, sobre todo, por la prescripción adquisitiva de diez años del art. 2459. Las contradicciones del art. 2461 con el art. 2449 debilitan la protección del legitimario. Muchas de estas cuestiones quedarán libradas a la doctrina y a la jurisprudencia.
} \\
\end{longtable}

\end{document}
